\documentclass[12pt,a4paper]{report}

\usepackage[T1]{fontenc}
\usepackage[utf8]{inputenc}
\usepackage[french]{babel}
\usepackage{geometry}
\usepackage{setspace}
\usepackage{lmodern}
\usepackage{tikz}
\usepackage{xcolor}
\usepackage{graphicx}
\usepackage{float}
\usepackage{array}
\usepackage{booktabs}
\usepackage{longtable}
\usepackage{caption}
\usepackage{listings}
\usepackage{hyperref}

\geometry{
    top=2.5cm,
    bottom=2.5cm,
    left=2.5cm,
    right=2.5cm
}

\setstretch{1.15}
\hypersetup{
    colorlinks=true,
    linkcolor=black,
    urlcolor=blue,
    citecolor=black
}

\definecolor{codebg}{RGB}{248,248,248}
\definecolor{codeframe}{RGB}{210,210,210}
\definecolor{codekeyword}{RGB}{0,75,135}
\definecolor{codestring}{RGB}{130,70,20}
\definecolor{codecomment}{RGB}{80,120,80}

\lstdefinestyle{pfeCode}{
    backgroundcolor=\color{codebg},
    frame=single,
    rulecolor=\color{codeframe},
    basicstyle=\ttfamily\small,
    keywordstyle=\color{codekeyword}\bfseries,
    stringstyle=\color{codestring},
    commentstyle=\color{codecomment}\itshape,
    breaklines=true,
    tabsize=2,
    showstringspaces=false,
    captionpos=b
}
\lstset{style=pfeCode}

% Figure utilisée pour les captures d'écran.
% Si le fichier image existe dans images/, il sera affiché.
% Sinon, un emplacement explicite reste visible dans le PDF.
\newcommand{\screenshotplaceholder}[3]{
\begin{figure}[H]
    \centering
    \IfFileExists{images/#1}{
        \includegraphics[width=0.92\textwidth]{images/#1}
    }{
        \fbox{
            \begin{minipage}[c][5cm][c]{0.86\textwidth}
                \centering
                \textbf{[Insérer ici la capture d'écran : #2]}\\[0.25cm]
                Fichier attendu : \texttt{images/\detokenize{#1}}
            \end{minipage}
        }
    }
    \caption{#3}
\end{figure}
}

% Figure utilisée pour les diagrammes Draw.io.
% Exporter le fichier .drawio en PNG dans diagrammes/ pour l'affichage final.
\newcommand{\diagramplaceholder}[3]{
\begin{figure}[H]
    \centering
    \IfFileExists{diagrammes/#1.png}{
        \includegraphics[width=0.95\textwidth]{diagrammes/#1.png}
    }{
        \fbox{
            \begin{minipage}[c][5.2cm][c]{0.88\textwidth}
                \centering
                \textbf{[Insérer ici le diagramme : #2]}\\[0.25cm]
                Source Draw.io : \texttt{diagrammes/\detokenize{#1}.drawio}\\
                Export PNG attendu : \texttt{diagrammes/\detokenize{#1}.png}
            \end{minipage}
        }
    }
    \caption{#3}
\end{figure}
}

\begin{document}

%==================================
% PAGE DE GARDE
%==================================

\begin{titlepage}
    \centering
    \vspace*{1cm}
    {\Large Ministère de l'Enseignement Supérieur et de la Recherche Scientifique}\\[0.3cm]
    {\Large Faculté des Sciences de Monastir}\\[1.5cm]
    {\Huge \textbf{Rapport de Projet de Fin d'Études}}\\[1.2cm]
    {\LARGE \textbf{Conception et réalisation d'une plateforme de gestion documentaire dynamique}}\\[1.5cm]
    \begin{flushleft}
        \large
        \textbf{Réalisé par :} \hspace{0.3cm} Oumar \\
        \textbf{Encadrant académique :} \hspace{0.3cm} À compléter \\
        \textbf{Encadrant professionnel :} \hspace{0.3cm} À compléter \\
        \textbf{Organisme d'accueil :} \hspace{0.3cm} À compléter
    \end{flushleft}
    \vfill
    {\large Année universitaire 2025--2026}
\end{titlepage}

\pagenumbering{roman}

\chapter*{Dédicaces}
\addcontentsline{toc}{chapter}{Dédicaces}
À compléter.

\chapter*{Remerciements}
\addcontentsline{toc}{chapter}{Remerciements}
À compléter.

%==================================
% TITRE PERSONNALISÉ
%==================================

\vspace*{0.5cm}

\noindent
\begin{tikzpicture}

% Rectangle noir
\fill[black] (0,0) rectangle (0.9,1.8);

% Ligne horizontale
\draw[line width=0.5pt] (1.2,0.9) -- (14.8,0.9);

% Titre
\node[
anchor=west,
font=\fontsize{24}{28}\bfseries
] at (1.2,0.25)
{TABLE DES MATIÈRES};

\end{tikzpicture}

\vspace{1.2cm}

\tableofcontents
\clearpage

\listoffigures
\addcontentsline{toc}{chapter}{Liste des figures}
\clearpage

\listoftables
\addcontentsline{toc}{chapter}{Liste des tableaux}
\clearpage

\pagenumbering{arabic}

%=================================================
% INTRODUCTION GENERALE
%=================================================

\chapter*{Introduction générale}
\addcontentsline{toc}{chapter}{Introduction générale}

Dans le contexte actuel de transformation numérique, la gestion de l'information et de la documentation constitue un enjeu stratégique pour toutes les organisations. La capacité à produire des documents fiables, cohérents et personnalisés, tout en garantissant la sécurité des accès, est devenue une nécessité opérationnelle. Pourtant, de nombreux processus administratifs reposent encore sur des manipulations manuelles de fichiers bureautiques, ce qui expose l'organisation à des erreurs de saisie, à des problèmes de versioning et à une hétérogénéité visuelle.

La problématique centrale de ce projet de fin d'études réside dans la conception et la réalisation d'une solution capable d'automatiser le cycle de vie documentaire. Il s'agit de passer d'une gestion statique, basée sur des fichiers isolés, à une gestion dynamique, où le document est le résultat d'une association contrôlée entre des modèles éditables, des familles de documents, des chartes graphiques et des données métier extraites directement du système d'information.

L'objectif de ce projet est donc de fournir une plateforme web complète et sécurisée permettant aux administrateurs de définir des structures documentaires flexibles et aux utilisateurs finaux de générer des documents personnalisés en quelques clics. La solution s'appuie sur une architecture moderne utilisant ASP.NET Core pour le backend, Angular pour le frontend et SQL Server pour la persistance des données.

Le présent rapport décrit la démarche de réalisation de cette plateforme, organisée en six chapitres :
\begin{itemize}
    \item le premier chapitre présente le cadre général du projet, l'organisme d'accueil et l'environnement technique ;
    \item le deuxième chapitre détaille l'analyse des besoins et la spécification fonctionnelle ;
    \item les chapitres trois à six décrivent la réalisation technique à travers les quatre sprints du projet, couvrant l'authentification, la gouvernance d'accès, l'édition documentaire, puis l'administration métier avec l'intégration finale.
\end{itemize}

L'ensemble de ce travail vise à démontrer comment une approche logicielle structurée peut répondre concrètement aux défis de la gestion documentaire moderne.



%=================================================
% CHAPITRES EXISTANTS
%=================================================

%=================================================
% CHAPITRE 1 : Cadre général du projet
%=================================================

\chapter{Cadre général du projet}

\section{Introduction}

Le point de départ de ce projet est un constat assez simple, mais très concret : dans beaucoup d'administrations scolaires ou universitaires, la gestion des ressources humaines ressemble encore à un puzzle de fichiers Excel et de documents Word copiés-collés. On passe un temps fou à ressaisir les mêmes noms, les mêmes grades et les mêmes dates dans des attestations de travail ou des contrats. 

L'idée de SIRH-Doc est née de cette frustration. L'objectif n'est pas seulement de stocker des noms dans une base de données, mais de créer une plateforme intelligente capable de transformer ces données en documents officiels, tout en gardant une trace impeccable de chaque décision administrative. Ce premier chapitre pose le décor : pourquoi ce projet, pour qui, et avec quels outils j'ai décidé de le construire.

\section{Présentation de l'organisme d'accueil}

Le projet s'inscrit dans le cadre d'un établissement d'enseignement supérieur. C'est un environnement où la gestion du personnel est particulièrement complexe : il faut jongler entre les enseignants permanents, les agents administratifs et une multitude de vacataires dont les contrats changent à chaque session.

Chaque jour, le service RH doit produire des dizaines de pièces : certificats d'enseignement, états de services, attestations de présence... La moindre erreur de frappe sur un grade ou une date peut invalider un document officiel. L'établissement disposait déjà de données informatisées, mais il manquait ce \og pont \fg{} automatisé entre la donnée brute et le document final prêt à imprimer.

\section{Présentation du projet}

\subsection{Contexte et Problématique}

Aujourd'hui, quand on veut créer une attestation pour un enseignant, on cherche souvent le dernier fichier Word qui a servi de modèle, on fait un \og Enregistrer sous \fg{}, et on modifie manuellement les informations. Cette méthode artisanale pose trois problèmes majeurs :
\begin{itemize}
    \item \textbf{Le risque d'erreur} : On oublie souvent de changer une date ou un titre caché en bas de page.
    \item \textbf{La perte d'historique} : Il est très difficile de savoir exactement ce qui a été délivré à qui il y a deux ans.
    \item \textbf{L'hétérogénéité} : Selon la personne qui rédige, le logo n'est pas à la même place ou la police change.
\end{itemize}

Ma problématique était donc la suivante : \textit{Comment centraliser la gestion du personnel et automatiser la production documentaire pour garantir à la fois un gain de temps et une fiabilité totale ?}

\subsection{La solution : SIRH-Doc}

La solution que j'ai conçue est une plateforme web intégrée. Elle ne se contente pas d'être un éditeur de texte ; c'est un véritable hub où l'administrateur définit des \og familles de documents \fg{}, configure des modèles dynamiques avec des variables (comme \texttt{\{\{nom\}\}} ou \texttt{\{\{grade\}\}}), et laisse le système fusionner tout cela avec les données réelles du personnel.

\section{Méthodologie de travail}

Même si j'ai mené ce développement seul, j'ai voulu suivre une démarche structurée en m'inspirant de la méthode **Agile Scrum**. J'ai découpé le projet en quatre sprints de trois à quatre semaines chacun. 

Cette approche m'a permis de ne pas me laisser submerger par la complexité de l'éditeur documentaire dès le début. J'ai d'abord sécurisé les fondations (Sprint 1), puis construit le moteur de données (Sprint 2), avant de m'attaquer au studio d'édition (Sprint 3) et à la production finale (Sprint 4).

\section{Environnement technique : Mes choix}

\subsection{Backend : .NET 8 et Dapper}
Pour le serveur, j'ai choisi **ASP.NET Core (Web API)**. C'est un framework que j'apprécie pour sa rigueur et sa performance. Pour l'accès aux données, j'ai préféré **Dapper** à Entity Framework. Pourquoi ? Parce que SIRH-Doc doit pouvoir exécuter des requêtes SQL très spécifiques et dynamiques selon les besoins des établissements, et Dapper me donne ce contrôle total sur le SQL tout en restant extrêmement rapide.

\subsection{Frontend : Angular 17}
Côté interface, **Angular** s'est imposé naturellement. Sa structure modulaire est parfaite pour une application d'administration lourde. J'ai utilisé la version 17 pour profiter des dernières optimisations de performance (comme le nouveau système de contrôle de flux).

\subsection{L'édition : Tiptap}
Pour la partie \og Studio \fg{}, j'ai intégré **Tiptap**. Contrairement aux éditeurs classiques qui sont parfois trop rigides, Tiptap est \og headless \fg{}. Cela m'a permis de construire une interface de création de modèles totalement personnalisée, tout en produisant un code HTML propre pour la génération des PDF.

\section{Conclusion}

Ce chapitre a permis de définir le \og pourquoi \fg{} du projet. SIRH-Doc n'est pas un simple exercice technique, c'est une réponse à un besoin de modernisation administrative. Maintenant que le cadre est posé, le chapitre suivant va détailler plus finement les fonctionnalités attendues et la manière dont j'ai modélisé le système.

%=================================================
% CHAPITRE 2
%=================================================

\chapter{Analyse et spécification des besoins}

\section{Introduction}

Après avoir présenté le contexte général du projet, ce chapitre se concentre sur l'analyse des besoins et la spécification fonctionnelle de la plateforme. Cette étape est essentielle, car elle permet de transformer une idée générale, à savoir automatiser et sécuriser la gestion documentaire, en exigences claires et exploitables par le développement.

L'analyse ne se limite pas à lister les fonctionnalités attendues. Elle cherche à comprendre les acteurs, leurs responsabilités, les informations manipulées, les contraintes de sécurité, les règles de gestion et les priorités du produit. Dans un système documentaire dynamique, une mauvaise compréhension des besoins peut rapidement provoquer des incohérences : un modèle mal relié à une famille, une charte appliquée au mauvais document, des bénéficiaires mal filtrés ou des droits accordés à un utilisateur non autorisé.

Ce chapitre présente donc les acteurs du système, clarifie le vocabulaire métier, détaille les besoins fonctionnels et non fonctionnels, expose le backlog produit et introduit la modélisation UML globale. Il constitue le lien entre le cadrage du chapitre précédent et la réalisation progressive décrite dans les sprints suivants.

\section{Identification des acteurs}

La plateforme distingue trois acteurs principaux. Cette séparation correspond à des responsabilités métier différentes et constitue la base du contrôle d'accès par rôles.

\subsection{Super administrateur}

Le super administrateur possède le niveau de responsabilité le plus élevé. Il intervient sur la configuration globale de la plateforme. Son rôle consiste à préparer les structures qui permettront ensuite aux administrateurs et aux utilisateurs de travailler dans un environnement cohérent.

Il peut gérer les familles de documents, configurer les vues de tables métier, vérifier les données disponibles, administrer les ressources structurantes et contrôler les paramètres partagés. Il est également responsable de la cohérence globale du référentiel documentaire. Par exemple, lorsqu'une nouvelle catégorie de documents doit être ajoutée, le super administrateur définit la famille correspondante, précise la source des bénéficiaires et configure les champs utiles.

\subsection{Administrateur}

L'administrateur est un utilisateur métier avancé. Il ne gère pas nécessairement toute la structure globale, mais il intervient sur les modèles et les ressources documentaires utilisées dans son périmètre. Il peut créer ou modifier des modèles, associer une charte graphique, organiser le contenu en en-tête, corps et pied de page, puis vérifier le rendu.

Son rôle demande une compréhension fonctionnelle du document attendu, mais ne doit pas nécessiter une intervention dans le code source. La plateforme doit donc lui fournir des interfaces suffisamment claires pour créer et maintenir des modèles sans dépendre d'un développeur.

\subsection{Utilisateur final}

L'utilisateur final exploite les modèles et les données disponibles pour produire ou prévisualiser un document. Il sélectionne une famille, choisit un modèle, recherche un bénéficiaire, applique éventuellement des filtres, puis consulte le rendu généré. Son périmètre doit être limité aux opérations nécessaires à son travail quotidien.

Cet acteur représente la finalité de la plateforme. Si l'administration prépare les ressources, l'utilisateur final valide la valeur réelle du système : produire rapidement un document fiable, correctement mis en forme et alimenté par les bonnes données.

\section{Clarification du vocabulaire métier}

La clarification du vocabulaire est importante, car plusieurs termes peuvent paraître proches alors qu'ils désignent des concepts différents. Cette distinction évite les ambiguïtés dans la conception et dans la rédaction des chapitres de réalisation.

\subsection{Famille de documents}

Une famille de documents représente une catégorie fonctionnelle. Elle regroupe des modèles qui partagent une logique métier commune. Elle peut définir la table des bénéficiaires, les colonnes d'identification, les champs d'affichage, les filtres disponibles et les variables exploitables dans les modèles.

Par exemple, une famille peut correspondre à un ensemble d'attestations, de décisions ou de documents administratifs destinés à un même type de bénéficiaire. Elle joue le rôle de pont entre le monde documentaire et les données métier.

\subsection{Modèle de document ou template}

Le modèle de document, également appelé template, décrit la structure d'un document. Il contient le contenu éditable, généralement divisé en en-tête, corps et pied de page. Il peut contenir du texte, des tableaux, des images, des variables et des paramètres de mise en page.

Le modèle n'est pas le document final. Il constitue une matrice réutilisable. Le document final est produit lorsque le modèle est combiné avec les données d'un bénéficiaire et une charte graphique.

\subsection{Charte graphique}

La charte graphique regroupe les règles visuelles appliquées aux documents : couleurs, typographies, identité de l'organisation, watermark, styles de titres et paramètres d'affichage. Elle permet de garantir une cohérence visuelle entre plusieurs modèles.

Centraliser la charte graphique évite de dupliquer les mêmes paramètres dans chaque modèle. Une modification visuelle peut ainsi être répercutée sur plusieurs documents sans reprendre individuellement tous les contenus.

\subsection{Configuration de vue de table}

La configuration de vue de table, représentée dans le code par \texttt{TableViewConfig}, décrit la manière dont une table métier est affichée et manipulée dans l'application. Elle précise le nom de la table, son libellé, les champs visibles, les champs éditables, les champs utilisés pour la prévisualisation et les éventuels paramètres de lookup.

Ce concept apporte une dimension configurable très importante. Au lieu de développer une interface spécifique pour chaque table métier, le système peut générer une vue exploitable à partir d'une configuration. Cela rend la plateforme plus souple et plus adaptée à des données évolutives.

\subsection{Instance de document}

Une instance de document correspond au résultat produit à partir d'un modèle, d'une charte graphique et des données d'un bénéficiaire. Elle représente le document personnalisé que l'utilisateur peut prévisualiser, imprimer ou exporter selon les fonctionnalités disponibles.

Cette distinction entre modèle et instance est essentielle. Le modèle est stable et réutilisable, tandis que l'instance dépend du bénéficiaire sélectionné et du contexte de génération.

\section{Analyse des besoins fonctionnels}

\subsection{Gestion des utilisateurs}

La plateforme doit permettre l'identification des utilisateurs à travers un mécanisme de connexion sécurisé. Chaque utilisateur possède un compte, une adresse électronique, un mot de passe haché, un rôle, éventuellement une organisation et des informations de profil. Le système doit permettre la connexion, l'inscription selon les règles prévues et la récupération du profil authentifié.

La gestion des utilisateurs constitue la première condition de sécurité. Sans identification fiable, il devient impossible de contrôler l'accès aux modèles, aux données métier et aux documents générés.

\subsection{Gestion des rôles et permissions}

Le système doit distinguer les rôles \texttt{supAdmin}, \texttt{admin} et \texttt{user}. Chaque rôle donne accès à un espace différent. Le super administrateur accède aux configurations globales, l'administrateur gère les ressources documentaires et l'utilisateur final exploite les modèles disponibles.

Le contrôle d'accès doit être appliqué côté frontend et côté backend. Les guards Angular empêchent l'accès aux routes non autorisées, tandis que les attributs d'autorisation et le middleware JWT protègent les endpoints. Cette double protection évite de dépendre uniquement de l'interface utilisateur.

\subsection{Gestion des familles de documents}

La plateforme doit permettre de créer, modifier, consulter et supprimer des familles de documents. Une famille doit pouvoir contenir des informations sur les bénéficiaires, les filtres, les colonnes d'affichage et les variables disponibles. Elle doit également servir de critère de regroupement pour les modèles.

Cette fonctionnalité est indispensable pour structurer le système documentaire. Sans famille, les modèles seraient isolés et difficiles à relier aux données métier.

\subsection{Gestion des modèles de documents}

Le système doit permettre la création et la modification de modèles de documents. Un modèle doit être associé à une famille et peut être associé à une charte graphique. Il doit contenir un en-tête, un corps et un pied de page. Il doit également conserver les paramètres de mise en page comme les marges, l'orientation et les règles d'affichage.

L'administrateur doit pouvoir éditer le contenu à travers un éditeur riche, sauvegarder le résultat et le prévisualiser. Cette fonctionnalité représente le coeur de la plateforme.

\subsection{Gestion des chartes graphiques}

La plateforme doit permettre de gérer des chartes graphiques. Une charte peut être liée à une organisation ou être définie comme charte par défaut. Elle contient une configuration visuelle au format structuré. Elle doit pouvoir être appliquée aux modèles afin de garantir une identité homogène.

\subsection{Gestion des configurations de vues de tables}

Le système doit permettre de gérer les configurations de vues de tables. Chaque configuration doit décrire une table métier, les champs affichables, les champs modifiables, les champs de prévisualisation et les réglages de lookup. L'utilisateur autorisé doit pouvoir consulter les lignes, créer un enregistrement, modifier un enregistrement ou supprimer une ligne selon les droits et la configuration.

Cette fonctionnalité rend l'application plus générique. Elle permet d'exploiter des données métier sans développer systématiquement un nouveau module.

\subsection{Édition et génération des documents}

La plateforme doit permettre d'éditer un modèle, d'y insérer du contenu riche et de générer une prévisualisation à partir d'un bénéficiaire. Les variables du modèle doivent être remplacées par les données disponibles. Le rendu doit appliquer les paramètres du modèle et de la charte graphique.

Cette fonctionnalité donne sa valeur principale au projet : transformer des données structurées en documents personnalisés et exploitables.

\subsection{Gestion des bénéficiaires et des données métier}

Le système doit permettre la recherche des bénéficiaires selon une famille. La recherche peut utiliser une table, une requête SQL configurée ou des filtres métier. Les données retournées doivent être limitées et adaptées au contexte de l'utilisateur.

La gestion des bénéficiaires est directement liée à la génération documentaire. Un document généré sans données fiables perd son intérêt ; il est donc essentiel que cette étape soit correctement conçue.

\section{Analyse des besoins non fonctionnels}

\subsection{Sécurité}

La sécurité constitue une exigence majeure. Les mots de passe doivent être hachés, les sessions doivent être représentées par des tokens JWT signés, les routes doivent être protégées et les endpoints sensibles doivent refuser les appels anonymes. Le système doit également éviter d'exposer des informations sensibles dans les réponses HTTP.

Les requêtes liées aux données métier doivent être contrôlées. Les configurations de vues de tables ne doivent pas devenir un moyen d'exécuter librement des requêtes dangereuses. Les champs manipulés doivent être limités aux champs configurés et autorisés.

\subsection{Performance}

L'application doit rester réactive même lorsque les données métier deviennent volumineuses. Les listes doivent être limitées, les ressources doivent être chargées progressivement et les appels API doivent éviter les transferts inutiles. Côté Angular, le chargement paresseux des composants contribue à réduire le poids initial de l'application.

\subsection{Ergonomie et expérience utilisateur}

L'ergonomie est importante, car la plateforme est destinée à des utilisateurs qui ne sont pas nécessairement développeurs. Les interfaces doivent guider l'utilisateur, réduire les manipulations inutiles et afficher des messages compréhensibles en cas d'erreur. L'éditeur documentaire doit offrir une expérience proche d'un outil bureautique tout en respectant les contraintes du système.

\subsection{Maintenabilité}

La maintenabilité repose sur la séparation des responsabilités. Le backend est organisé en contrôleurs, services, repositories, DTOs et entités. Le frontend est organisé en composants, services, guards, interceptors et modèles. Cette structure permet d'ajouter de nouvelles fonctionnalités sans modifier toute l'application.

Les configurations JSON, comme les chartes graphiques ou les vues de tables, apportent de la flexibilité. Elles doivent toutefois être normalisées pour éviter l'apparition de données incohérentes.

\section{Backlog du produit}

\subsection{User stories principales}

Le backlog produit rassemble les besoins exprimés sous forme de user stories. Il représente une vision globale du produit avant son découpage en sprints.

\begin{longtable}{p{0.16\textwidth} p{0.42\textwidth} p{0.30\textwidth}}
\caption{Backlog produit global}\\
\toprule
\textbf{ID} & \textbf{User story} & \textbf{Valeur attendue}\\
\midrule
\endfirsthead
\toprule
\textbf{ID} & \textbf{User story} & \textbf{Valeur attendue}\\
\midrule
\endhead
US-01 & En tant qu'utilisateur, je peux me connecter de manière sécurisée. & Protéger l'accès à la plateforme.\\
\midrule
US-02 & En tant qu'utilisateur connecté, je suis redirigé vers l'espace correspondant à mon rôle. & Améliorer l'expérience et appliquer la gouvernance d'accès.\\
\midrule
US-03 & En tant que super administrateur, je peux gérer les familles de documents. & Structurer les catégories documentaires.\\
\midrule
US-04 & En tant qu'administrateur, je peux créer et modifier des modèles. & Produire des matrices documentaires réutilisables.\\
\midrule
US-05 & En tant qu'administrateur, je peux définir une charte graphique. & Garantir une identité visuelle homogène.\\
\midrule
US-06 & En tant que super administrateur, je peux configurer les vues de tables. & Administrer les données métier de manière configurable.\\
\midrule
US-07 & En tant qu'utilisateur, je peux rechercher un bénéficiaire. & Personnaliser le document avec les bonnes données.\\
\midrule
US-08 & En tant qu'utilisateur, je peux prévisualiser le document généré. & Vérifier le résultat avant exploitation.\\
\midrule
US-09 & En tant que système, je dois refuser les accès non autorisés. & Assurer la confidentialité et l'intégrité.\\
\bottomrule
\end{longtable}

\subsection{Priorisation des besoins}

La priorisation suit une logique de dépendance. L'authentification et les rôles sont prioritaires, car ils conditionnent l'accès au reste de l'application. L'édition documentaire vient ensuite, car elle représente la valeur centrale du produit. Les modules d'administration métier complètent le système en rendant les modèles exploitables dans un contexte réel. Enfin, l'intégration et les tests consolident la solution.

\subsection{Critères d'acceptation}

Les critères d'acceptation permettent de vérifier objectivement si une fonctionnalité est terminée. Ils sont formulés de manière observable.

\begin{longtable}{p{0.28\textwidth} p{0.58\textwidth}}
\caption{Critères d'acceptation principaux}\\
\toprule
\textbf{Fonctionnalité} & \textbf{Critères d'acceptation}\\
\midrule
\endfirsthead
\toprule
\textbf{Fonctionnalité} & \textbf{Critères d'acceptation}\\
\midrule
\endhead
Authentification & Un utilisateur valide reçoit un JWT ; un utilisateur invalide reçoit une erreur ; le mot de passe n'est jamais retourné.\\
\midrule
RBAC & Chaque rôle accède uniquement à son espace ; une route interdite provoque une redirection ; un endpoint protégé refuse les appels anonymes.\\
\midrule
Modèles & Un modèle peut être créé, modifié, sauvegardé et rechargé sans perte de contenu.\\
\midrule
Chartes graphiques & Une charte peut être sauvegardée et appliquée au rendu documentaire.\\
\midrule
TableViewConfig & Une vue de table peut être configurée ; seuls les champs visibles et éditables déclarés sont utilisés.\\
\midrule
Génération & Un document peut être prévisualisé à partir d'un modèle et d'un bénéficiaire.\\
\bottomrule
\end{longtable}

\section{Planification des sprints}

\subsection{Découpage temporel}

Le projet est découpé en quatre sprints. Ce découpage permet de passer progressivement de l'infrastructure vers la valeur métier, puis de regrouper la consolidation finale avec l'administration opérationnelle. Chaque sprint produit un livrable vérifiable et prépare le sprint suivant.

\subsection{Objectifs par sprint}

Le premier sprint concerne l'authentification et le socle applicatif. Le deuxième sprint met en place la gestion des rôles. Le troisième sprint développe l'édition documentaire. Le quatrième sprint ajoute l'administration des modules métier et prend en charge l'intégration, les tests et la livraison.

\subsection{Livrables attendus}

Les livrables attendus sont les fichiers sources backend et frontend, les interfaces fonctionnelles, les endpoints testables, les diagrammes UML, les captures d'écran, la documentation technique et le présent rapport.

\section{Diagrammes UML globaux}

\subsection{Diagramme de cas d'utilisation général}

\diagramplaceholder{global_cas_utilisation}{Diagramme de cas d'utilisation global}{Vue globale des acteurs et fonctionnalités du système}

Le diagramme global de cas d'utilisation présente les fonctionnalités principales accessibles selon les acteurs. Il montre que l'authentification est un point d'entrée commun, puis que les responsabilités se séparent selon les rôles. Le super administrateur se concentre sur la configuration, l'administrateur sur les ressources documentaires et l'utilisateur final sur l'exploitation des documents.

\subsection{Diagramme de classes global}

\diagramplaceholder{global_classes}{Diagramme de classes global}{Vue globale des entités, services, DTOs et repositories}

Le diagramme de classes global donne une vue d'ensemble de l'architecture. Il relie les entités \texttt{User}, \texttt{Family}, \texttt{Template}, \texttt{GraphicCharter} et \texttt{TableViewConfig} aux services et repositories. Il montre également le rôle des DTOs dans les échanges entre frontend et backend.

\subsection{Diagramme de séquence global}

\diagramplaceholder{global_sequences}{Diagramme de séquence global}{Flux global Angular, API, services, repository et SQL Server}

Le diagramme de séquence global décrit un parcours complet : authentification, chargement de l'état, sélection d'un modèle, recherche d'un bénéficiaire, génération de la prévisualisation et retour du document rendu. Il permet de visualiser la collaboration entre Angular, l'API ASP.NET Core, les services backend, les repositories et SQL Server.

\section{Architecture globale du système}

L'architecture globale repose sur une séparation claire entre le frontend, le backend et la base de données. Angular prend en charge les interfaces, la navigation, les guards, les interceptors et les services clients. ASP.NET Core expose les endpoints REST, applique les règles de sécurité, orchestre la logique métier et interagit avec SQL Server. La base de données conserve les données persistantes et les informations métier.

Cette architecture présente plusieurs avantages. Elle sépare l'expérience utilisateur de la logique serveur, facilite les tests, permet l'évolution indépendante des couches et offre une base solide pour de futures extensions, comme l'export PDF avancé, la signature électronique ou l'historisation des documents générés.

\section{Analyse technique globale du projet}

L'analyse du projet SIRH-Doc montre une architecture organisée autour de trois responsabilités principales. Le backend ASP.NET Core expose les contrats REST, applique les règles de sécurité et centralise l'accès aux données. Le frontend Angular porte l'expérience utilisateur, les écrans d'administration, l'éditeur documentaire et les services clients. La base SQL Server conserve à la fois les données de configuration documentaire et les données métier utilisées pour alimenter les documents.

Côté backend, la séparation entre contrôleurs, services, repositories et DTOs rend le code lisible et maintenable. Les contrôleurs restent proches du protocole HTTP, tandis que les services portent l'orchestration applicative. Les repositories isolent les requêtes SQL et facilitent l'évolution des accès aux données. Cette organisation n'est pas un DDD strict au sens tactique complet, mais elle reprend une logique de séparation par responsabilités qui convient au périmètre du projet.

La gestion multi-tenant constitue un point structurant. Le tenant courant permet d'associer les traitements à une organisation et à une base de données cible. L'ajout de l'année universitaire active renforce cette logique, car plusieurs tables métier doivent être filtrées selon la période académique. Ce mécanisme améliore la traçabilité des dossiers du personnel et évite de mélanger des données appartenant à des exercices universitaires différents.

Le système de génération documentaire repose sur les familles, les modèles, les chartes graphiques, les variables simples, les listes et les tableaux dynamiques. Cette conception rend le moteur documentaire configurable : un administrateur peut faire évoluer un modèle sans modifier le code source. Les modules dynamiques, notamment les vues de tables, complètent cette logique en permettant d'exposer des données métier sans créer une interface spécifique pour chaque table.

La sécurité repose sur l'authentification JWT, le contrôle des rôles, les guards Angular, les middlewares backend et la validation des requêtes. Les points sensibles concernent surtout l'exécution de requêtes SQL configurables et l'accès aux données multi-tenant. Ces risques sont limités par l'autorisation des requêtes de lecture uniquement, la normalisation des paramètres, le filtrage par organisation et l'ajout de contrôles liés à l'année universitaire.

Sur le plan UX/UI, l'application adopte une organisation par espaces : super administrateur, administrateur et utilisateur. Cette séparation clarifie les responsabilités, mais impose de garder une cohérence forte dans les formulaires, les messages d'erreur et les écrans de configuration. Les améliorations appliquées pendant la consolidation ont principalement visé la fiabilité de l'éditeur, la génération des tableaux dynamiques, la gestion des images et le chargement complet des données nécessaires.

En matière de maintenabilité, le projet présente une base saine : services spécialisés, modèles typés côté frontend, repositories centralisés côté backend et composants dédiés pour les fonctions principales. Les axes d'amélioration à poursuivre concernent l'augmentation des tests automatisés, la documentation des conventions SQL configurables, la surveillance des performances sur les gros volumes et la mise en place progressive d'un audit plus détaillé des actions sensibles.

\section{Conclusion}

Ce chapitre a permis de formaliser les besoins de la plateforme. Les acteurs ont été identifiés, les concepts métier clarifiés, les fonctionnalités décrites et les contraintes non fonctionnelles précisées. Le backlog produit et les critères d'acceptation donnent une base de suivi pour la réalisation.

La modélisation UML globale complète cette analyse en offrant une vision structurée du système. Les chapitres suivants décrivent maintenant la réalisation progressive de la plateforme à travers les quatre sprints définis.


%=================================================
% CHAPITRES REDIGES A PARTIR DU SPRINT 1
%=================================================

%=================================================
% CHAPITRE 3 : SPRINT 1
%=================================================

\chapter{Sprint 1 : Socle technique, Multi-Tenancy et isolation temporelle}

\section{Introduction}

Quand j'ai commencé à poser les premières briques du projet, ma priorité n'était pas de construire de jolis écrans, mais de m'assurer que les données soient parfaitement isolées. Dans un contexte universitaire, on manipule des informations sensibles (diplômes, contrats, salaires) et il est impensable que les données d'une faculté puissent être visibles par une autre, ou que les contrats de l'année passée se mélangent avec ceux de l'année en cours.

Ce premier sprint a donc servi à bâtir les fondations invisibles mais indispensables de SIRH-Doc : l'architecture multi-tenant et la gestion du contexte temporel par année universitaire.

\section{Objectifs et Backlog du sprint}

L'idée était d'arriver, à la fin de cette itération, à une application capable de reconnaître qui se connecte, pour quelle organisation il travaille, et sur quelle période académique il souhaite agir.

\begin{itemize}
    \item \textbf{Mise en place du Multi-Tenancy} : Capacité du backend à aiguiller les requêtes vers la bonne base de données.
    \item \textbf{Sécurité JWT} : Authentification robuste avec des jetons qui portent l'identité et le périmètre de l'utilisateur.
    \item \textbf{Contexte d'Année Universitaire} : Création d'un filtre global pour cloisonner les données RH par session.
\end{itemize}

\section{Choix d'architecture et défis rencontrés}

\subsection{La résolution dynamique du Tenant}

Au départ, j'avais envisagé une gestion classique via session utilisateur côté serveur, mais cette approche devenait vite compliquée avec les changements d'organisation. J'ai finalement opté pour un middleware de \textit{Tenant Resolution}. À chaque appel HTTP, le système regarde le jeton de l'utilisateur et configure instantanément la connexion SQL via une \texttt{ConnectionFactory}. C'est transparent pour le reste du code, ce qui m'a évité de passer l'identifiant de l'organisation dans chaque fonction.

\subsection{L'année universitaire comme contexte global}

C'est sans doute le point qui m'a demandé le plus de réflexion. On ne peut pas simplement mettre une date dans une table. Un vacataire peut avoir un contrat en 2024-2025 et un autre en 2025-2026. J'ai donc choisi de traiter l'année universitaire comme un \og état global \fg{} de l'application. Une fois sélectionnée côté frontend, elle est injectée dans toutes les requêtes SQL via Dapper pour filtrer les parcours et les positions administratives.

\section{Analyse et Conception}

\subsection{Diagramme de cas d'utilisation}

Le diagramme suivant montre comment le point d'entrée du système est verrouillé. Avant de pouvoir faire quoi que ce soit, l'utilisateur doit franchir l'étape de l'authentification et du choix de son contexte de travail (organisation et année).

\diagramplaceholder{sprint1_cas_utilisation}{Cas d'utilisation du Sprint 1}{Initialisation du contexte et sécurité}

\subsection{Diagramme de classes (Entités persistantes)}

Conformément à mon choix d'isoler les responsabilités, les classes persistantes de ce sprint se concentrent sur la gestion des accès et des structures de base. On y retrouve l'utilisateur, son organisation de rattachement et le référentiel des années académiques.

\diagramplaceholder{sprint1_classes}{Diagramme de classes - Sprint 1}{Entités de base pour la sécurité et le contexte}

\subsection{Diagrammes de séquence : Flux techniques du socle}

Les flux suivants illustrent les deux mécanismes fondamentaux de ce sprint : la sécurisation des accès et l'isolation automatique des données.

\subsubsection{Authentification et génération de jeton (JWT)}

Le flux de connexion illustre bien la complexité cachée derrière un simple bouton \og Valider \fg{}. Le backend ne se contente pas de vérifier un mot de passe ; il construit un jeton JWT contenant le périmètre de l'utilisateur (rôle et organisation) pour sécuriser les appels futurs.

\diagramplaceholder{sprint1_sequence_auth}{Séquence d'authentification}{Processus de login et génération du périmètre utilisateur}

\subsubsection{Résolution dynamique du Tenant (Middleware)}

Ce diagramme détaille le fonctionnement du middleware .NET. À chaque requête, le système intercepte le jeton JWT et les en-têtes HTTP pour identifier l'établissement et l'année universitaire. Cette information est injectée dans un contexte \textit{Scoped} pour filtrer automatiquement toutes les requêtes SQL.

\diagramplaceholder{sprint1_sequence_middleware}{Séquence de résolution de Tenant}{Interception et injection du contexte organisationnel}

\section{Réalisation technique}

\subsection{Backend : Sécurisation et Dapper}

Pour le stockage des mots de passe, je n'ai pris aucun risque : j'ai utilisé \textbf{BCrypt}. Côté accès aux données, \textbf{Dapper} s'est avéré être un excellent choix pour sa rapidité, surtout quand il s'agit d'injecter dynamiquement des clauses \texttt{WHERE} liées au TenantId ou à l'année universitaire.

\subsection{Frontend : Le Guard de sélection}

Côté Angular, j'ai mis en place un \textit{Guard} spécifique. Si un administrateur tente d'accéder à la gestion du personnel sans avoir choisi une année universitaire active, il est automatiquement renvoyé vers un écran de sélection. Cela garantit qu'aucune donnée n'est créée \og dans le vide \fg{}.

\section{Conclusion du sprint}

À la fin de ce sprint, j'avais un socle solide. L'application savait qui était connecté et dans quel cadre temporel il travaillait. C'était l'étape nécessaire avant de pouvoir attaquer la gestion réelle des dossiers RH dans le sprint suivant.

%=================================================
% CHAPITRE 4
%=================================================

\chapter{Sprint 2 : Gestion des rôles et gouvernance d'accès}

\section{Introduction}

Après la mise en place de l'authentification, le deuxième sprint a été consacré à la gouvernance d'accès. Dans une plateforme documentaire, tous les utilisateurs ne doivent pas disposer du même périmètre fonctionnel. Un super administrateur peut gérer la structure globale de l'application, un administrateur peut intervenir sur les modèles et les données de son périmètre, tandis qu'un utilisateur final doit accéder principalement aux documents et aux données qui lui sont autorisés.

Ce sprint introduit donc le contrôle d'accès basé sur les rôles, également appelé RBAC pour \textit{Role-Based Access Control}. Le principe consiste à associer chaque utilisateur à un rôle, puis à autoriser ou refuser les fonctionnalités selon ce rôle. Cette logique doit être appliquée à deux niveaux : côté frontend, pour orienter l'utilisateur vers les bonnes interfaces, et côté backend, pour empêcher les accès directs aux endpoints protégés.

La difficulté de ce sprint réside dans la cohérence entre l'expérience utilisateur et la sécurité réelle. Masquer un menu dans Angular améliore l'ergonomie, mais ne suffit pas à sécuriser l'application. Les contrôles doivent donc être dupliqués de manière complémentaire : les guards protègent la navigation côté client, tandis que les attributs d'autorisation et la validation du JWT protègent les ressources côté serveur.

\section{Objectifs du sprint}

Les objectifs du sprint visent à transformer l'authentification en véritable mécanisme d'autorisation. Le système ne doit plus seulement savoir qui est connecté, mais aussi ce que cet utilisateur a le droit de faire.

Les objectifs fonctionnels et techniques sont les suivants :

\begin{itemize}
    \item définir les rôles \texttt{supAdmin}, \texttt{admin} et \texttt{user} ;
    \item intégrer le rôle dans l'entité utilisateur, dans le token JWT et dans l'objet utilisateur retourné au frontend ;
    \item mettre en place une redirection automatique vers l'espace correspondant au rôle ;
    \item protéger les routes Angular à l'aide d'un \texttt{roleGuard} ;
    \item sécuriser les endpoints backend avec \texttt{[Authorize]} et les claims de rôle ;
    \item préparer les interfaces distinctes : espace super administrateur, espace administrateur et espace utilisateur ;
    \item valider les comportements en cas d'accès non autorisé.
\end{itemize}

Ces objectifs assurent une gouvernance minimale mais robuste, indispensable avant d'exposer les fonctionnalités de gestion documentaire.

\section{Backlog du sprint}

\begin{longtable}{p{0.17\textwidth} p{0.38\textwidth} p{0.35\textwidth}}
\caption{Backlog du sprint 2}\\
\toprule
\textbf{User story} & \textbf{Description} & \textbf{Tâches principales}\\
\midrule
\endfirsthead
\toprule
\textbf{User story} & \textbf{Description} & \textbf{Tâches principales}\\
\midrule
\endhead
US2.1 & En tant que super administrateur, je peux accéder à l'espace global d'administration. & Définir le rôle \texttt{supAdmin}, créer la route Angular dédiée, autoriser uniquement ce rôle.\\
\midrule
US2.2 & En tant qu'administrateur, je peux accéder à l'espace de gestion métier. & Définir le rôle \texttt{admin}, créer la route \texttt{/admin}, contrôler les accès aux modèles et documents.\\
\midrule
US2.3 & En tant qu'utilisateur final, je peux accéder uniquement à mon espace documentaire. & Définir le rôle \texttt{user}, créer la route \texttt{/user}, limiter les fonctionnalités visibles.\\
\midrule
US2.4 & En tant qu'utilisateur connecté, je suis automatiquement redirigé vers mon espace. & Implémenter la redirection dans \texttt{loginGuard} et dans la réponse d'authentification.\\
\midrule
US2.5 & En tant que système, je dois refuser les accès non autorisés. & Tester les accès directs aux routes, vérifier les réponses \texttt{401} et \texttt{403}, protéger les endpoints sensibles.\\
\bottomrule
\end{longtable}

\subsection{Définition des rôles}

Trois rôles ont été retenus. Le rôle \texttt{supAdmin} représente le niveau d'administration le plus élevé. Il permet de gérer les configurations globales, les familles, les vues de tables et les paramètres structurants. Le rôle \texttt{admin} correspond à un administrateur métier chargé de gérer les modèles, les chartes graphiques et les documents associés à son périmètre. Le rôle \texttt{user} correspond à l'utilisateur final, qui consulte les données autorisées, sélectionne un modèle et génère ou prévisualise un document.

Cette répartition répond au besoin de séparation des responsabilités. Elle évite de donner à un utilisateur final des droits de modification sur les modèles et limite les erreurs de manipulation sur les données structurantes.

\subsection{Protection des routes et des endpoints}

La protection des routes Angular est assurée par \texttt{roleGuard}. Ce guard vérifie d'abord que l'utilisateur est authentifié, puis compare son rôle avec la liste des rôles autorisés dans la configuration de la route. En cas d'incompatibilité, l'utilisateur est redirigé vers son propre espace.

Côté backend, la protection repose sur le middleware JWT et les attributs \texttt{[Authorize]}. Les endpoints qui manipulent des ressources sensibles ne doivent pas accepter les appels anonymes. Cette protection serveur reste indispensable car un utilisateur peut appeler directement l'API sans passer par l'interface Angular.

\subsection{Gestion des utilisateurs et profils}

La gestion des profils s'appuie sur les champs \texttt{Role}, \texttt{OrganizationId}, \texttt{Profile}, \texttt{ProfileDetail}, \texttt{AccessAllYears} et \texttt{AccessYearList}. Ces informations permettent d'aller au-delà du simple rôle technique. Elles préparent les restrictions métier par organisation, par profil fonctionnel ou par période.

\section{Analyse et conception}

\subsection{Diagramme de cas d'utilisation}

\diagramplaceholder{sprint2_cas_utilisation}{Diagramme de cas d'utilisation du sprint 2}{Cas d'utilisation de la gouvernance d'accès}

Le diagramme de cas d'utilisation du sprint 2 distingue trois acteurs : super administrateur, administrateur et utilisateur final. Chaque acteur hérite du comportement commun de l'utilisateur authentifié, mais dispose d'un périmètre différent.

Le super administrateur peut accéder à l'espace de configuration globale, consulter l'état de l'application et préparer les données de référence. L'administrateur peut gérer les contenus documentaires et les ressources liées à son organisation. L'utilisateur final peut consulter les modèles disponibles et interagir avec les documents générés. Le cas \og Accéder à une route protégée \fg{} inclut systématiquement la vérification du token et du rôle.

\subsection{Diagramme de classes}

\diagramplaceholder{sprint2_classes}{Diagramme de classes du sprint 2}{Classes et services impliqués dans le RBAC}

Le diagramme de classes met en évidence les éléments techniques du contrôle d'accès : \texttt{User}, \texttt{AuthUserResponse}, \texttt{AuthService}, \texttt{roleGuard}, \texttt{loginGuard}, \texttt{AuthInterceptor} et les composants de pages. Le rôle est présent dans l'entité persistée, dans le token JWT et dans l'objet utilisateur stocké côté Angular.

Cette continuité est importante. Si le rôle était seulement stocké côté frontend, il pourrait être altéré. S'il était seulement en base, chaque navigation Angular nécessiterait une requête supplémentaire. Le compromis retenu consiste à placer le rôle dans le JWT signé et à le conserver dans l'objet utilisateur du frontend pour l'ergonomie.

\subsection{Diagrammes de séquence}

\diagramplaceholder{sprint2_sequences}{Diagramme de séquence du sprint 2}{Séquence d'accès à une route protégée}

Le diagramme de séquence illustre le scénario d'accès à une route protégée. Lorsqu'un utilisateur tente d'ouvrir \texttt{/super-admin}, \texttt{/admin} ou \texttt{/user}, Angular appelle \texttt{roleGuard}. Celui-ci vérifie l'existence d'une session valide via \texttt{AuthService}. Si l'utilisateur n'est pas authentifié, il est redirigé vers \texttt{/login}. S'il est authentifié, son rôle est comparé au tableau \texttt{roles} défini dans la route.

Lorsque la route est autorisée, le composant correspondant est chargé de manière paresseuse. Lorsqu'elle ne l'est pas, l'utilisateur est redirigé vers l'espace adapté à son rôle. Cette logique améliore l'expérience utilisateur en évitant l'affichage d'une erreur brute.

\subsubsection*{Description des interactions système}

La gouvernance d'accès combine plusieurs interactions. Le backend émet un token contenant le rôle. Le frontend stocke ce rôle dans l'objet utilisateur courant. Les routes Angular définissent les rôles autorisés. Le guard applique la décision d'accès avant le chargement du composant. Enfin, l'interceptor ajoute le token aux appels API pour que le backend puisse appliquer sa propre validation.

\subsubsection*{Analyse technique}

Le \texttt{roleGuard} est implémenté sous forme de guard fonctionnel Angular. Cette approche correspond au style moderne d'Angular et évite de créer une classe de guard plus lourde. Le routage utilise le chargement dynamique des composants, ce qui réduit le poids initial de l'application.

\begin{lstlisting}[caption={Configuration des routes protégées par rôle}]
{
  path: "admin",
  loadComponent: () =>
    import("./features/editor/pages/admin/admin.component")
      .then((m) => m.AdminComponent),
  canActivate: [roleGuard],
  data: { page: "admin", roles: ["admin"] },
}
\end{lstlisting}

\subsubsection*{Architecture des composants concernés}

L'architecture RBAC s'articule autour de trois couches. La couche routeur décide si une page peut être activée. La couche service conserve l'état d'authentification et expose l'utilisateur courant. La couche API valide les tokens et protège les endpoints. Ce découpage permet une sécurité progressive : l'interface guide l'utilisateur, tandis que l'API garantit la protection effective des données.

\section{Réalisation}

\subsection{Implémentation du système RBAC}

Le RBAC a été implémenté en capitalisant sur le token JWT généré au sprint précédent. Le claim \texttt{ClaimTypes.Role} est ajouté au token lors de la connexion. Le frontend reçoit également le rôle dans l'objet \texttt{AuthUserResponse}. Cette information permet de rediriger l'utilisateur immédiatement après l'authentification.

\begin{lstlisting}[caption={Redirection selon le rôle après authentification}]
private static string GetRoleHome(string? role) => role switch
{
    "supAdmin" => "/superAdmin.html",
    "admin" => "/admin.html",
    _ => "/user.html"
};
\end{lstlisting}

Même si cette redirection est utile, la sécurité réelle reste portée par le contrôle d'accès. Le rôle dans la réponse ne doit jamais être considéré comme une preuve autonome. La preuve d'identité reste le token signé.

\subsection{Mise en place des autorisations backend}

Les endpoints sensibles utilisent l'attribut \texttt{[Authorize]}. Cet attribut déclenche la validation du token avant l'exécution de l'action du contrôleur. Par exemple, la récupération des templates exige une session valide.

\begin{lstlisting}[caption={Protection d'un endpoint de templates}]
[HttpGet]
[Authorize]
public async Task<ActionResult<IEnumerable<object>>> GetAll()
{
    return Ok(await _service.GetTemplatesAsync(CurrentUser()));
}
\end{lstlisting}

Le contrôleur peut ensuite récupérer les informations de l'utilisateur courant à partir des claims. Cette méthode permet notamment de filtrer les ressources selon l'organisation ou le rôle.

\subsection{Gestion des guards côté frontend}

Le guard de rôle vérifie la session et compare le rôle de l'utilisateur avec les rôles autorisés. En cas de refus, il redirige vers l'espace approprié.

\begin{lstlisting}[caption={Vérification du rôle dans le guard Angular}]
const allowedRoles: string[] = route.data["roles"] ?? [];
if (!allowedRoles.length) return true;

const user = authService.getCurrentUser();
if (user && allowedRoles.includes(user.role)) {
  return true;
}

const role = user?.role;
if (role === "supAdmin") router.navigate(["/super-admin"]);
else if (role === "admin") router.navigate(["/admin"]);
else router.navigate(["/user"]);
\end{lstlisting}

Cette logique évite qu'un utilisateur soit bloqué sans indication. Elle renforce également l'ergonomie, car l'utilisateur retrouve automatiquement son espace naturel.

\screenshotplaceholder{sprint2_routes.png}{Routes protégées selon les rôles}{Navigation protégée par rôle dans l'application Angular}
\screenshotplaceholder{sprint2_admin.png}{Interface administrateur}{Espace administrateur accessible après contrôle RBAC}

\subsection{Sécurisation des accès par profil}

Au-delà du rôle, le profil utilisateur peut contenir des informations complémentaires. Les champs d'organisation et de profil sont transmis dans le JWT et l'objet utilisateur. Cette conception prépare les filtrages métier, par exemple l'affichage de modèles liés à une organisation ou la limitation des données par année.

Les contrôleurs peuvent construire un objet utilisateur courant à partir des claims :

\begin{lstlisting}[caption={Construction de l'utilisateur courant à partir des claims}]
return new Dictionary<string, object?>
{
    ["id"] = User.FindFirstValue(ClaimTypes.NameIdentifier),
    ["name"] = User.FindFirstValue(ClaimTypes.Name),
    ["email"] = User.FindFirstValue(ClaimTypes.Email),
    ["role"] = User.FindFirstValue(ClaimTypes.Role),
    ["organizationId"] = User.FindFirstValue("organizationId"),
    ["profile"] = ""
};
\end{lstlisting}

Cette méthode sert de pont entre la couche sécurité et la couche métier. Elle permet au service documentaire de charger un état adapté au contexte de l'utilisateur.

\section{Tests et validation}

Les tests du sprint 2 ont vérifié la cohérence entre les routes Angular, les rôles utilisateurs et les endpoints backend.

\begin{longtable}{p{0.25\textwidth} p{0.43\textwidth} p{0.22\textwidth}}
\caption{Scénarios de test du sprint 2}\\
\toprule
\textbf{Scénario} & \textbf{Description} & \textbf{Résultat attendu}\\
\midrule
\endfirsthead
\toprule
\textbf{Scénario} & \textbf{Description} & \textbf{Résultat attendu}\\
\midrule
\endhead
Accès super administrateur & Connexion avec un compte \texttt{supAdmin} puis accès à \texttt{/super-admin}. & Accès autorisé.\\
\midrule
Accès admin interdit au super espace & Connexion avec un compte \texttt{admin} puis tentative d'accès à \texttt{/super-admin}. & Redirection vers \texttt{/admin}.\\
\midrule
Accès utilisateur final & Connexion avec un compte \texttt{user} puis accès à \texttt{/user}. & Accès autorisé.\\
\midrule
Route protégée sans session & Accès direct à \texttt{/admin} sans token. & Redirection vers \texttt{/login}.\\
\midrule
Endpoint protégé sans token & Appel direct à \texttt{/api/templates}. & Réponse \texttt{401 Unauthorized}.\\
\bottomrule
\end{longtable}

Les validations techniques ont confirmé que les composants protégés n'étaient pas chargés lorsque le rôle ne correspondait pas. Les tests de sécurité ont montré qu'un accès direct à un endpoint protégé sans token était rejeté par le backend, indépendamment des contrôles effectués côté client.

\section{Conclusion}

Le deuxième sprint a transformé le socle d'authentification en un système de gouvernance d'accès cohérent. Les rôles sont désormais intégrés au token, aux objets utilisateur, aux routes Angular et aux contrôles backend. La plateforme distingue clairement les espaces \texttt{supAdmin}, \texttt{admin} et \texttt{user}.

Cette étape prépare le développement du coeur métier. Le sprint suivant pourra se concentrer sur l'édition documentaire avec Tiptap, en s'appuyant sur une navigation sécurisée et sur une identité utilisateur déjà maîtrisée.

%=================================================
% CHAPITRE 5
%=================================================

\chapter{Sprint 3 : Module coeur d'édition documentaire}

\section{Introduction}

Le troisième sprint représente le coeur fonctionnel du projet. Après avoir sécurisé l'accès à l'application et mis en place la gouvernance par rôles, il devient possible de développer le module d'édition documentaire. Ce module permet de créer, modifier, sauvegarder et prévisualiser des modèles de documents qui seront ensuite utilisés pour produire des documents personnalisés à partir des données métier.

La plateforme ne se limite pas à stocker des fichiers statiques. Elle propose une logique documentaire dynamique fondée sur des modèles composés d'un en-tête, d'un corps et d'un pied de page. Ces modèles peuvent contenir du texte riche, des tableaux, des images, des variables et des paramètres de mise en page. Le choix de l'éditeur Tiptap répond à ce besoin, car il offre une base extensible pour construire un éditeur WYSIWYG moderne au sein d'une application Angular.

Ce sprint a donc pour objectif de livrer une première version exploitable de l'éditeur documentaire. L'utilisateur doit pouvoir saisir un contenu structuré, l'enregistrer dans un modèle, le recharger ultérieurement et obtenir une prévisualisation cohérente avec les paramètres de mise en page.

\section{Objectifs du sprint}

Les objectifs du sprint couvrent à la fois l'expérience utilisateur et la robustesse technique de la gestion du contenu documentaire :

\begin{itemize}
    \item intégrer Tiptap comme éditeur de texte riche dans l'application Angular ;
    \item organiser le modèle documentaire en zones distinctes : en-tête, corps et pied de page ;
    \item permettre la création et la modification du contenu HTML associé à un modèle ;
    \item sauvegarder les modèles via les services Angular et les endpoints REST ;
    \item charger les modèles depuis l'état applicatif et les afficher dans l'éditeur ;
    \item appliquer les paramètres de mise en page tels que l'orientation, les marges et les distances d'en-tête et de pied de page ;
    \item produire une prévisualisation documentaire à partir d'un modèle et d'un bénéficiaire ;
    \item préparer l'intégration des chartes graphiques et des variables métier.
\end{itemize}

\section{Backlog du sprint}

\begin{longtable}{p{0.17\textwidth} p{0.38\textwidth} p{0.35\textwidth}}
\caption{Backlog du sprint 3}\\
\toprule
\textbf{User story} & \textbf{Description} & \textbf{Tâches principales}\\
\midrule
\endfirsthead
\toprule
\textbf{User story} & \textbf{Description} & \textbf{Tâches principales}\\
\midrule
\endhead
US3.1 & En tant qu'administrateur, je peux éditer le contenu d'un modèle de document. & Intégrer Tiptap, créer les zones d'édition, gérer le contenu HTML.\\
\midrule
US3.2 & En tant qu'administrateur, je peux sauvegarder les modifications d'un modèle. & Normaliser le modèle, appeler \texttt{TemplateService}, persister via \texttt{PUT /api/templates/\{id\}}.\\
\midrule
US3.3 & En tant qu'utilisateur, je peux prévisualiser le document généré. & Implémenter \texttt{DocumentRenderService}, appeler l'endpoint \texttt{/api/preview}, remplacer les variables.\\
\midrule
US3.4 & En tant qu'administrateur, je peux gérer l'en-tête et le pied de page du modèle. & Ajouter les attributs \texttt{hasHeader}, \texttt{hasFooter}, \texttt{headerDisplay}, \texttt{footerDisplay}.\\
\midrule
US3.5 & En tant que système, je dois conserver une structure cohérente du modèle. & Définir \texttt{TemplateRecord}, normaliser les données, valider les champs essentiels.\\
\bottomrule
\end{longtable}

\subsection{Intégration de l'éditeur riche}

L'intégration de l'éditeur riche s'appuie sur Tiptap, un éditeur basé sur ProseMirror. Ce choix se justifie par sa modularité, son support des extensions et sa capacité à produire un contenu HTML structuré. L'objectif n'est pas seulement de disposer d'une zone de texte, mais de fournir un environnement d'édition compatible avec les besoins documentaires : mise en forme, tableaux, images, styles, alignements et variables.

\subsection{Création et modification du contenu documentaire}

La création du contenu documentaire se fait au niveau du modèle. Le modèle contient les champs \texttt{header}, \texttt{body} et \texttt{footer}, qui correspondent respectivement à l'en-tête, au corps et au pied de page. Cette séparation rend la génération plus maîtrisable, car chaque zone peut avoir ses propres règles d'affichage.

\subsection{Sauvegarde et chargement des contenus}

La sauvegarde est assurée par \texttt{TemplateService}. Le service normalise le modèle, met à jour l'état local, puis transmet les modifications au backend. Cette stratégie améliore la réactivité de l'interface tout en conservant une persistance côté serveur.

\subsection{Prévisualisation du document}

La prévisualisation permet à l'utilisateur de vérifier le rendu avant l'utilisation finale du document. Elle combine le contenu du modèle, les données du bénéficiaire, les paramètres de mise en page et les éventuelles chartes graphiques. Elle constitue une étape essentielle pour garantir la qualité documentaire.

\section{Analyse et conception}

\subsection{Diagramme de cas d'utilisation}

\diagramplaceholder{sprint3_cas_utilisation}{Diagramme de cas d'utilisation du sprint 3}{Cas d'utilisation du module d'édition documentaire}

Le diagramme de cas d'utilisation présente deux acteurs principaux : l'administrateur et l'utilisateur final. L'administrateur gère les modèles, édite leur contenu et sauvegarde les modifications. L'utilisateur final sélectionne un modèle disponible, choisit un bénéficiaire et prévisualise le document généré.

Le cas d'utilisation \og Prévisualiser un document \fg{} inclut la récupération du modèle, la récupération des données du bénéficiaire et le rendu HTML final. Le cas \og Sauvegarder un modèle \fg{} inclut la validation du contenu et la persistance via l'API.

\subsection{Diagramme de classes}

\diagramplaceholder{sprint3_classes}{Diagramme de classes du sprint 3}{Classes du module d'édition et de rendu documentaire}

Le diagramme de classes montre les relations entre \texttt{TemplateRecord}, \texttt{TemplateService}, \texttt{DocumentRenderService}, \texttt{DocumentDataService}, \texttt{GraphicCharterService} et les composants Angular d'administration et d'utilisation. Le modèle documentaire est l'entité centrale. Il référence une famille, une organisation éventuelle et une charte graphique.

La classe backend \texttt{Template} contient des champs proches du modèle frontend : identifiant, famille, organisation, charte graphique, nom, en-tête, corps, pied de page, marges et orientation. Cette correspondance facilite la sérialisation et limite les transformations inutiles.

\subsection{Diagrammes de séquence}

\diagramplaceholder{sprint3_sequences}{Diagramme de séquence du sprint 3}{Séquence d'édition, sauvegarde et prévisualisation}

Le diagramme de séquence du sprint décrit deux scénarios complémentaires. Dans le premier, l'administrateur sélectionne un modèle, modifie son contenu dans Tiptap, puis déclenche la sauvegarde. Le composant Angular appelle \texttt{TemplateService}, qui met à jour l'état local et envoie une requête \texttt{PUT} à l'API. Le backend transmet la demande à \texttt{EditorService}, puis au repository.

Dans le deuxième scénario, l'utilisateur demande une prévisualisation. Le frontend récupère le modèle et les données du bénéficiaire, puis demande au service de rendu de construire le HTML final. Lorsque le rendu nécessite des données serveur, l'endpoint \texttt{/api/preview} remplace les variables simples du modèle par les valeurs du bénéficiaire.

\subsubsection*{Description des interactions système}

L'interaction d'édition est centrée sur le modèle. L'éditeur modifie le contenu, le composant prépare un objet \texttt{TemplateRecord}, puis le service garantit sa normalisation. L'interaction de prévisualisation est plus riche : elle combine modèle, bénéficiaire, données métier, charte graphique et règles de pagination.

\subsubsection*{Analyse technique}

Le contenu documentaire est stocké sous forme HTML. Cette décision facilite l'intégration avec Tiptap et permet une restitution directe dans le navigateur. Elle exige toutefois des précautions : le contenu doit être maîtrisé, les champs variables doivent suivre une convention claire, et les traitements serveur doivent éviter l'exécution de SQL ou de HTML non contrôlé.

Le backend expose un endpoint \texttt{/api/preview}. Celui-ci récupère le modèle, assemble l'en-tête, le corps et le pied de page, puis remplace les variables de type \texttt{\{\{champ\}\}} ou \texttt{\{champ\}} lorsque les données du bénéficiaire sont disponibles.

\subsubsection*{Architecture des composants concernés}

Le module s'appuie sur une architecture frontend orientée services. Le composant d'administration gère l'interaction utilisateur. \texttt{TemplateService} persiste les modèles. \texttt{DocumentDataService} récupère les données liées aux familles. \texttt{DocumentRenderService} produit le HTML final. \texttt{GraphicCharterService} fournit les variables visuelles de la charte.

\section{Réalisation}

\subsection{Intégration de l'éditeur Tiptap}

L'intégration de Tiptap a consisté à installer les extensions nécessaires et à créer une interface d'édition riche. Les extensions utilisées couvrent les besoins classiques d'un document professionnel : paragraphes, titres, listes, tableaux, images, couleurs, polices et alignements. L'éditeur est intégré dans les pages d'administration afin de modifier les modèles.

\screenshotplaceholder{sprint3_tiptap.png}{Éditeur Tiptap intégré}{Interface d'édition riche d'un modèle documentaire}

Le contenu de chaque zone est synchronisé avec le modèle sélectionné. Lorsqu'un modèle change, l'éditeur recharge les valeurs \texttt{header}, \texttt{body} et \texttt{footer}. Lors de la sauvegarde, ces contenus sont réintégrés dans l'objet \texttt{TemplateRecord}.

\subsection{Gestion du contenu des modèles de documents}

Le modèle frontend est représenté par l'interface \texttt{TemplateRecord}. Il contient l'identifiant, la famille, l'organisation, la charte graphique, les zones HTML et les paramètres de rendu.

\begin{lstlisting}[caption={Structure simplifiée d'un modèle documentaire côté Angular}]
export interface TemplateRecord {
  id: string;
  familyId?: string;
  organizationId?: string | null;
  graphicCharterId?: string | null;
  header?: string;
  body?: string;
  footer?: string;
  hasHeader?: boolean;
  hasFooter?: boolean;
  filterProfile: TemplateFilterProfileEntry[];
  sectionDirections: TemplateSectionDirections;
}
\end{lstlisting}

Le service de sauvegarde met à jour l'état local avant d'appeler l'API. Cela permet à l'interface de rester réactive tout en garantissant la persistance.

\begin{lstlisting}[caption={Sauvegarde d'un modèle par TemplateService}]
async saveTemplate(template: TemplateRecord): Promise<TemplateRecord> {
  const normalized = normalizeTemplateRecord(template);
  const state = this.state.getState();
  const index = state.templates.findIndex(item => item.id === normalized.id);
  index >= 0
    ? state.templates.splice(index, 1, normalized)
    : state.templates.push(normalized);
  this.state.replaceState(state);
  await firstValueFrom(
    this.api.put(`templates/${encodeURIComponent(normalized.id)}`, normalized)
  );
  return normalized;
}
\end{lstlisting}

\subsection{Application des en-têtes, corps et pieds de page}

Le découpage du document en trois zones permet de répondre à des besoins professionnels : logo ou informations institutionnelles en en-tête, contenu principal dans le corps, signatures ou références en pied de page. Les champs \texttt{hasHeader} et \texttt{hasFooter} permettent d'activer ou de désactiver ces zones selon le type de document.

L'application prend également en compte les directions de section. Cette fonctionnalité est utile dans un contexte multilingue, par exemple pour gérer des contenus en français et en arabe dans un même système documentaire.

\subsection{Prévisualisation et rendu des documents}

La prévisualisation est prise en charge par \texttt{DocumentRenderService}. Ce service construit un rendu HTML à partir du modèle et des données du bénéficiaire. Il applique les variables CSS issues de la charte graphique, les marges de page, les distances d'en-tête et de pied de page, ainsi que les règles de pagination.

\begin{lstlisting}[caption={Principe de remplacement des variables dans l'endpoint de prévisualisation}]
html = Regex.Replace(html,
    "\\{\\{\\s*(\\w+)\\s*\\}\\}|\\{\\s*(\\w+)\\s*\\}",
    match =>
    {
        var key = match.Groups[1].Success
            ? match.Groups[1].Value
            : match.Groups[2].Value;
        return rowValues.TryGetValue(key, out var val) && val != null
            ? System.Net.WebUtility.HtmlEncode(val)
            : match.Value;
    });
\end{lstlisting}

\screenshotplaceholder{sprint3_preview.png}{Prévisualisation documentaire}{Prévisualisation d'un document généré à partir d'un modèle}

Cette implémentation donne une première base de génération dynamique. Elle pourra ensuite être enrichie avec l'export PDF, la gestion avancée des blocs conditionnels et l'intégration de variables complexes.

\section{Tests et validation}

\begin{longtable}{p{0.25\textwidth} p{0.43\textwidth} p{0.22\textwidth}}
\caption{Scénarios de test du sprint 3}\\
\toprule
\textbf{Scénario} & \textbf{Description} & \textbf{Résultat attendu}\\
\midrule
\endfirsthead
\toprule
\textbf{Scénario} & \textbf{Description} & \textbf{Résultat attendu}\\
\midrule
\endhead
Chargement d'un modèle & Sélectionner un modèle existant depuis l'interface admin. & Les zones d'édition affichent le contenu enregistré.\\
\midrule
Modification du contenu & Modifier le corps du document puis sauvegarder. & Le contenu est persisté et reste disponible après rechargement.\\
\midrule
Gestion de l'en-tête & Activer ou désactiver l'en-tête. & Le rendu reflète l'état du champ \texttt{hasHeader}.\\
\midrule
Prévisualisation & Sélectionner un bénéficiaire et lancer la prévisualisation. & Les variables du modèle sont remplacées par les données disponibles.\\
\midrule
Validation API & Appeler \texttt{PUT /api/templates/\{id\}} avec un token valide. & Réponse positive et modèle mis à jour.\\
\bottomrule
\end{longtable}

Les tests fonctionnels ont confirmé que l'édition, la sauvegarde et le rechargement du contenu fonctionnent de manière cohérente. Les tests d'intégration ont vérifié que le frontend et le backend manipulent une structure compatible du modèle. Les validations techniques ont porté sur la conservation des champs structurants et sur le rendu des zones HTML dans la prévisualisation.

\section{Conclusion}

Le troisième sprint a permis de livrer le coeur métier de la plateforme : l'édition documentaire dynamique. Grâce à Tiptap, aux services Angular et aux endpoints de templates, l'application peut désormais gérer des modèles de documents riches et prévisualiser un rendu personnalisé.

Ce sprint prépare directement l'administration métier du sprint suivant. Les modèles ne prennent tout leur sens que lorsqu'ils sont reliés à des familles de documents, à des chartes graphiques, à des configurations de vues de tables et à des bénéficiaires. Le sprint 4 sera donc consacré à la structuration de ces modules d'administration.

%=================================================
% CHAPITRE 6
%=================================================

\chapter{Sprint 4 : Administration des modules métier}

\section{Introduction}

Le quatrième sprint a pour objectif de compléter le coeur documentaire par les modules métier nécessaires à son exploitation. Un éditeur de modèles ne suffit pas à lui seul à produire des documents fiables. Il faut également définir les familles de documents, les chartes graphiques, les configurations de vues de tables, les bénéficiaires et les filtres permettant de sélectionner les données métier utilisées dans la génération.

Ce sprint représente donc une étape structurante. Il transforme l'application en une véritable plateforme d'administration documentaire. Les familles permettent de classer les modèles et de décrire les sources de données associées. Les modèles définissent la structure documentaire. Les chartes graphiques assurent l'identité visuelle. Les configurations de vues de tables, représentées par l'entité \texttt{TableViewConfig}, permettent d'exposer des données métier sous forme de vues administrables. Enfin, les bénéficiaires et les filtres permettent de sélectionner les données injectées dans les documents.

La complexité de ce sprint vient de l'interdépendance entre ces modules. Un modèle appartient à une famille, peut être lié à une organisation et peut utiliser une charte graphique. Une famille peut référencer une table de bénéficiaires ou une requête SQL. Une configuration de vue de table décrit les champs visibles, modifiables et utilisés pour la prévisualisation. L'enjeu est donc de proposer une administration cohérente, tout en conservant une architecture maintenable.

\section{Objectifs du sprint}

Les objectifs du sprint sont les suivants :

\begin{itemize}
    \item implémenter le CRUD des familles de documents ;
    \item implémenter le CRUD des modèles de documents ;
    \item implémenter le CRUD des chartes graphiques ;
    \item implémenter le CRUD des configurations de vues de tables \texttt{TableViewConfig} ;
    \item permettre la consultation, la création, la modification et la suppression des enregistrements métier à travers les vues de tables ;
    \item associer les modèles aux familles et aux chartes graphiques ;
    \item configurer les filtres et les bénéficiaires utilisés pour la génération documentaire ;
    \item valider les données envoyées aux endpoints et normaliser les structures JSON échangées.
\end{itemize}

\section{Backlog du sprint}

\begin{longtable}{p{0.17\textwidth} p{0.38\textwidth} p{0.35\textwidth}}
\caption{Backlog du sprint 4}\\
\toprule
\textbf{User story} & \textbf{Description} & \textbf{Tâches principales}\\
\midrule
\endfirsthead
\toprule
\textbf{User story} & \textbf{Description} & \textbf{Tâches principales}\\
\midrule
\endhead
US4.1 & En tant que super administrateur, je peux gérer les familles de documents. & Créer les endpoints \texttt{/api/families}, définir \texttt{Family}, gérer les bénéficiaires et filtres.\\
\midrule
US4.2 & En tant qu'administrateur, je peux gérer les modèles d'une famille. & Utiliser \texttt{/api/templates}, associer chaque modèle à une famille et une charte.\\
\midrule
US4.3 & En tant qu'administrateur, je peux définir des chartes graphiques. & Créer \texttt{GraphicCharter}, gérer les couleurs, polices, watermark et identité visuelle.\\
\midrule
US4.4 & En tant que super administrateur, je peux configurer les vues de tables métier. & Implémenter \texttt{TableViewConfig}, champs visibles, champs éditables, champs de prévisualisation et paramètres de lookup.\\
\midrule
US4.5 & En tant qu'utilisateur, je peux sélectionner un bénéficiaire à partir des données métier. & Implémenter \texttt{/api/search-beneficiaries}, appliquer les filtres et retourner les lignes disponibles.\\
\bottomrule
\end{longtable}

\subsection{Gestion des familles de documents}

Une famille de documents définit un regroupement fonctionnel. Elle précise le nom, la description, le mode de bénéficiaire, la table source, les colonnes d'affichage, la colonne de liaison, les filtres disponibles et les classes de variables. Elle sert donc de point d'entrée entre la logique documentaire et les données métier.

\subsection{Gestion des modèles de documents}

Le modèle est la structure éditable utilisée pour produire un document. Il est associé à une famille à travers \texttt{familyId}. Cette relation garantit que les variables et les bénéficiaires disponibles correspondent au type de document traité.

\subsection{Gestion des chartes graphiques}

La charte graphique permet de centraliser les paramètres visuels : couleurs, typographies, identité, watermark et règles d'affichage. Elle évite de dupliquer les styles dans chaque modèle et permet de faire évoluer l'identité visuelle d'une organisation de manière contrôlée.

\subsection{Gestion des configurations de vues de tables}

La configuration de vue de table est représentée par l'entité \texttt{TableViewConfig}. Elle décrit comment une table métier doit être présentée et éditée depuis l'application. Elle contient notamment le nom de la table, le libellé, les champs visibles, les champs éditables, les champs de prévisualisation et les paramètres de lookup.

Ce module est important, car il permet d'administrer des données métier sans développer une page spécifique pour chaque table. Il apporte donc une dimension configurable à la plateforme.

\subsection{Gestion des bénéficiaires et filtres métier}

Les bénéficiaires représentent les enregistrements métier utilisés pour personnaliser un document. Selon la famille, ils peuvent provenir d'une table SQL, d'une organisation ou d'une requête spécifique. Les filtres permettent de réduire la liste des bénéficiaires selon des critères configurés, par exemple une année, une direction, un service ou un statut.

\section{Analyse et conception}

\subsection{Diagramme de cas d'utilisation}

\diagramplaceholder{sprint4_cas_utilisation}{Diagramme de cas d'utilisation du sprint 4}{Cas d'utilisation d'administration des modules métier}

Le diagramme de cas d'utilisation du sprint 4 distingue principalement le super administrateur et l'administrateur. Le super administrateur gère les structures globales : familles, vues de tables et sources de données. L'administrateur gère les modèles et chartes graphiques dans son périmètre. L'utilisateur final intervient indirectement à travers la sélection des bénéficiaires et l'utilisation des données configurées.

Le cas d'utilisation \og Configurer une vue de table \fg{} inclut la sélection de la table, la définition des champs visibles, la définition des champs éditables et la configuration des champs de lookup. Le cas \og Gérer les bénéficiaires \fg{} dépend de la configuration des familles et des vues de tables.

\subsection{Diagramme de classes}

\diagramplaceholder{sprint4_classes}{Diagramme de classes du sprint 4}{Classes métier administrées dans le sprint 4}

Le diagramme de classes met en évidence les entités principales : \texttt{Family}, \texttt{Template}, \texttt{GraphicCharter} et \texttt{TableViewConfig}. Une famille possède plusieurs modèles. Un modèle peut référencer une charte graphique. Une configuration de vue de table référence une table métier et décrit les champs exploitables par l'interface.

Les services \texttt{FamilyService}, \texttt{TemplateService}, \texttt{TableViewService} et \texttt{GraphicCharterService} côté Angular communiquent avec \texttt{ApiService}. Côté backend, les contrôleurs dédiés délèguent les opérations à \texttt{IEditorService}, qui s'appuie sur \texttt{IEditorRepository}.

\subsection{Diagrammes de séquence}

\diagramplaceholder{sprint4_sequences}{Diagramme de séquence du sprint 4}{Séquence CRUD d'une configuration métier}

Le diagramme de séquence illustre le CRUD d'une configuration de vue de table. Le super administrateur crée une nouvelle configuration dans l'interface. Le composant Angular appelle \texttt{TableViewService.saveTableView}. Le service normalise les données, met à jour l'état local, puis envoie la requête au backend. Le contrôleur reçoit le JSON, appelle \texttt{EditorService}, puis le repository exécute l'insertion ou la mise à jour en base.

Un deuxième flux concerne la consultation des lignes métier. L'utilisateur choisit une vue de table. L'application appelle \texttt{/api/table-view/rows}. Le backend résout la configuration, vérifie les champs autorisés et construit une requête SQL limitée aux champs déclarés dans la configuration.

\subsubsection*{Description des interactions système}

Les interactions du sprint 4 sont principalement des interactions CRUD. Elles doivent toutefois rester cohérentes avec l'état global de l'application. Lorsqu'une famille est supprimée, les modèles associés doivent être pris en compte. Lorsqu'une charte graphique est modifiée, les modèles qui y font référence doivent afficher le nouveau rendu. Lorsqu'une vue de table est modifiée, les écrans de données métier doivent refléter les nouveaux champs visibles et éditables.

\subsubsection*{Analyse technique}

Le backend utilise des structures JSON pour gérer certaines entités configurables. Cette approche apporte de la souplesse, car les champs comme \texttt{filterCatalog}, \texttt{classes}, \texttt{fieldSettings} ou \texttt{pageMargins} peuvent évoluer sans modifier immédiatement toutes les classes métier. En contrepartie, elle impose une normalisation rigoureuse côté service et frontend.

Pour les vues de tables, la sécurité est particulièrement importante. Le backend ne doit pas exécuter librement des champs ou tables fournis par le client. La configuration doit être résolue, les champs doivent être comparés au schéma autorisé, et les requêtes doivent limiter les colonnes aux champs déclarés.

\subsubsection*{Architecture des composants concernés}

L'architecture du sprint suit une logique verticale par module. Chaque ressource dispose d'un modèle, d'un service Angular, d'un endpoint REST, d'un service backend et d'une méthode repository. Cette organisation facilite la compréhension et permet d'ajouter de nouvelles ressources métier selon le même modèle.

\section{Réalisation}

\subsection{CRUD des familles de documents}

Le CRUD des familles est exposé par \texttt{FamiliesController}. Les endpoints couvrent la liste, la consultation par identifiant, la création, la modification et la suppression.

\begin{lstlisting}[caption={Contrôleur REST pour les familles de documents}]
[ApiController]
[Route("api/families")]
public class FamiliesController : ControllerBase
{
    [HttpGet]
    public async Task<ActionResult<IEnumerable<object>>> GetAll()
    {
        return Ok(await _service.GetFamiliesAsync());
    }

    [HttpPut("{id}")]
    public async Task<ActionResult<object>> Update(string id, [FromBody] JsonObject family)
    {
        family["id"] = id;
        return Ok(await _service.UpsertFamilyAsync(family));
    }
}
\end{lstlisting}

\screenshotplaceholder{sprint4_families.png}{Gestion des familles de documents}{Interface de gestion des familles de documents}

\subsection{CRUD des modèles de documents}

Les modèles sont gérés à travers \texttt{TemplatesController}. Contrairement à certaines ressources publiques, les endpoints de templates sont protégés par \texttt{[Authorize]}. Cette protection est logique, car les modèles contiennent des contenus documentaires et des paramètres métier.

\begin{lstlisting}[caption={Endpoint de mise à jour d'un modèle}]
[HttpPut("{id}")]
[Authorize]
public async Task<ActionResult<object>> Update(string id, [FromBody] JsonObject template)
{
    template["id"] = id;
    return Ok(await _service.UpsertTemplateAsync(template));
}
\end{lstlisting}

\screenshotplaceholder{sprint4_templates.png}{Gestion des modèles de documents}{Interface de gestion des modèles associés aux familles}

\subsection{CRUD des chartes graphiques}

Les chartes graphiques sont gérées par \texttt{GraphicChartersController}. L'entité \texttt{GraphicCharter} contient un identifiant, une organisation éventuelle, un nom, une description, un indicateur de charte par défaut et un objet \texttt{Config}. Ce dernier regroupe les paramètres visuels.

\begin{lstlisting}[caption={Structure de l'entité GraphicCharter}]
public class GraphicCharter
{
    public string Id { get; set; } = string.Empty;
    public string? OrganizationId { get; set; }
    public string Name { get; set; } = string.Empty;
    public string? Description { get; set; }
    public bool IsDefault { get; set; }
    public JsonObject Config { get; set; } = new();
}
\end{lstlisting}

\screenshotplaceholder{sprint4_charters.png}{Gestion des chartes graphiques}{Interface de paramétrage d'une charte graphique}

\subsection{CRUD des configurations de vues de tables}

Le module \texttt{TableViewConfig} permet de configurer dynamiquement l'affichage et l'édition de tables métier. Il est exposé par deux familles d'endpoints : les endpoints REST classiques \texttt{/api/table-view-configs} et les endpoints opérationnels \texttt{/api/table-view-config}, \texttt{/api/table-view/rows}, \texttt{/api/table-view/record} et \texttt{/api/table-view/lookup-options}.

\begin{lstlisting}[caption={Entité TableViewConfig}]
public class TableViewConfig
{
    public string Id { get; set; } = string.Empty;
    public string TableName { get; set; } = string.Empty;
    public string Label { get; set; } = string.Empty;
    public JsonArray VisibleFields { get; set; } = new();
    public JsonArray EditableFields { get; set; } = new();
    public JsonArray PreviewFields { get; set; } = new();
    public JsonObject FieldSettings { get; set; } = new();
}
\end{lstlisting}

Côté Angular, \texttt{TableViewService} encapsule la sauvegarde, la suppression, la consultation des lignes, la création d'un enregistrement, la modification d'un enregistrement et la récupération des options de lookup.

\begin{lstlisting}[caption={Sauvegarde d'une configuration de vue de table}]
async saveTableView(view: TableViewConfig): Promise<TableViewConfig> {
  const normalized = normalizeTableViewRecord(view);
  const state = this.state.getState();
  const index = state.tableViews.findIndex(item => item.id === normalized.id);
  index >= 0
    ? state.tableViews.splice(index, 1, normalized)
    : state.tableViews.push(normalized);
  this.state.replaceState(state);
  await firstValueFrom(this.api.post('table-view-config',
    { tableView: normalized }));
  return normalized;
}
\end{lstlisting}

\screenshotplaceholder{sprint4_tableviews.png}{Gestion des configurations de vues de tables}{Interface de configuration et d'exploitation des vues de tables}

\subsection{Association famille, modèle et charte graphique}

L'association entre famille, modèle et charte graphique est assurée par les champs \texttt{familyId} et \texttt{graphicCharterId} du modèle. Cette association permet de déterminer les variables disponibles, les bénéficiaires pertinents et le style visuel à appliquer lors de la prévisualisation.

Dans l'interface, l'administrateur sélectionne une famille, puis crée ou modifie les modèles qui lui appartiennent. Il peut ensuite choisir la charte graphique à associer au modèle. Cette organisation évite les incohérences, par exemple l'utilisation d'un modèle avec des variables incompatibles avec la famille sélectionnée.

\subsection{Configuration des données métier et des filtres}

La recherche des bénéficiaires est assurée par l'endpoint \texttt{/api/search-beneficiaries}. Celui-ci reçoit l'identifiant de la famille, l'organisation, les filtres, une chaîne de recherche et une limite. Le backend charge la famille, lit la configuration du bénéficiaire, construit les paramètres et exécute la requête adaptée.

\begin{lstlisting}[caption={Validation minimale de la recherche de bénéficiaires}]
if (!request.TryGetProperty("familyId", out var famIdElem) ||
    famIdElem.ValueKind != JsonValueKind.String)
{
    return BadRequest(new EditorApiResponse(false,
        Error: "familyId is required"));
}
\end{lstlisting}

Cette fonctionnalité relie directement l'administration métier à la génération documentaire. Les bénéficiaires sélectionnés deviennent les sources de données utilisées pour remplacer les variables dans les modèles.

\section{Tests et validation}

\begin{longtable}{p{0.25\textwidth} p{0.43\textwidth} p{0.22\textwidth}}
\caption{Scénarios de test du sprint 4}\\
\toprule
\textbf{Scénario} & \textbf{Description} & \textbf{Résultat attendu}\\
\midrule
\endfirsthead
\toprule
\textbf{Scénario} & \textbf{Description} & \textbf{Résultat attendu}\\
\midrule
\endhead
Création famille & Créer une famille avec table bénéficiaire et colonnes d'affichage. & Famille enregistrée et visible dans la liste.\\
\midrule
Association modèle & Créer un modèle associé à une famille. & Le modèle apparaît dans la famille correspondante.\\
\midrule
Charte graphique & Modifier les couleurs et la typographie d'une charte. & Le rendu du document applique les paramètres.\\
\midrule
Configuration vue table & Définir champs visibles, éditables et preview. & La vue affiche uniquement les champs configurés.\\
\midrule
CRUD enregistrement métier & Créer, modifier et supprimer une ligne via \texttt{TableViewService}. & Les opérations sont répercutées en base.\\
\midrule
Recherche bénéficiaire & Rechercher un bénéficiaire avec filtre. & La liste retournée respecte la famille et les filtres.\\
\bottomrule
\end{longtable}

Les tests fonctionnels ont validé les opérations CRUD principales. Les tests d'intégration ont vérifié la cohérence entre les services Angular, les endpoints REST et la base SQL Server. Les tests de sécurité ont porté sur la limitation des champs manipulés par les vues de tables et sur le refus des requêtes incomplètes, notamment l'absence de \texttt{familyId} dans la recherche des bénéficiaires.

\section{Conclusion}

Le sprint 4 a transformé la plateforme en un système documentaire administrable. Les familles, modèles, chartes graphiques, configurations de vues de tables et bénéficiaires sont désormais structurés et exploitables. L'ajout explicite de \texttt{TableViewConfig} constitue un apport important, car il permet d'administrer les données métier de manière configurable sans développer une interface spécifique pour chaque table.

À l'issue de ce sprint, les principales fonctionnalités métier sont présentes. Le sprint suivant peut donc se concentrer sur l'intégration finale, les tests, l'optimisation, la documentation et la préparation de la livraison.

%=================================================
% CHAPITRE 7
%=================================================

\chapter{Sprint 5 : Intégration, qualité et livraison}

\section{Introduction}

Le cinquième sprint correspond à la phase de consolidation du projet. Après avoir développé le socle d'authentification, la gouvernance d'accès, l'édition documentaire et les modules métier, il devient nécessaire de vérifier la cohérence globale de l'application. Cette étape vise à transformer un ensemble de fonctionnalités développées progressivement en une solution stable, testable et présentable dans le cadre d'un projet de fin d'études.

L'intégration finale couvre plusieurs dimensions : intégration frontend/backend, validation des parcours utilisateurs, tests fonctionnels, tests d'intégration, optimisation des performances, renforcement de la sécurité et préparation de la documentation. Elle ne consiste donc pas seulement à corriger des anomalies. Elle permet également d'améliorer l'expérience utilisateur, de fiabiliser les échanges API et de préparer la démonstration finale.

Ce sprint a également une importance académique. Il permet de montrer que le projet ne se limite pas à une implémentation technique, mais suit une démarche de qualité logicielle. Les tests, la documentation, les captures d'écran, les diagrammes et la préparation de la livraison contribuent à rendre le travail évaluable, reproductible et maintenable.

\section{Objectifs du sprint}

Les objectifs du sprint 5 sont les suivants :

\begin{itemize}
    \item intégrer et valider les flux complets entre Angular et ASP.NET Core ;
    \item vérifier les parcours principaux des trois rôles : super administrateur, administrateur et utilisateur final ;
    \item tester les opérations CRUD sur les familles, modèles, chartes graphiques et configurations de vues de tables ;
    \item valider le processus de génération documentaire dynamique ;
    \item améliorer la gestion des erreurs et l'affichage des messages utilisateur ;
    \item renforcer les contrôles de sécurité sur les endpoints sensibles ;
    \item optimiser le chargement des ressources et la réactivité des interfaces ;
    \item préparer la documentation technique, la documentation utilisateur et la démonstration PFE.
\end{itemize}

\section{Backlog du sprint}

\begin{longtable}{p{0.17\textwidth} p{0.38\textwidth} p{0.35\textwidth}}
\caption{Backlog du sprint 5}\\
\toprule
\textbf{User story} & \textbf{Description} & \textbf{Tâches principales}\\
\midrule
\endfirsthead
\toprule
\textbf{User story} & \textbf{Description} & \textbf{Tâches principales}\\
\midrule
\endhead
US5.1 & En tant qu'utilisateur final, je peux utiliser l'application sans blocage majeur. & Tester connexion, navigation, sélection famille, sélection modèle, prévisualisation et génération.\\
\midrule
US5.2 & En tant qu'administrateur, je peux gérer les ressources métier de bout en bout. & Tester CRUD modèles, chartes, familles et vues de tables.\\
\midrule
US5.3 & En tant que système, je dois afficher des erreurs compréhensibles. & Améliorer l'interceptor d'erreur, valider les réponses API, afficher les messages utiles.\\
\midrule
US5.4 & En tant qu'évaluateur PFE, je peux comprendre et reproduire l'exécution. & Rédiger documentation, préparer captures, produire diagrammes et décrire l'architecture.\\
\midrule
US5.5 & En tant que développeur, je peux livrer une version stable. & Exécuter build, tests prioritaires, revue des endpoints, vérification sécurité.\\
\bottomrule
\end{longtable}

\subsection{Intégration frontend/backend}

L'intégration frontend/backend consiste à vérifier que les contrats HTTP sont cohérents. Les services Angular doivent appeler les bons endpoints, envoyer les données sous la forme attendue et interpréter correctement les réponses. Côté backend, les contrôleurs doivent retourner des statuts HTTP cohérents et des messages d'erreur exploitables.

\subsection{Tests fonctionnels et tests d'intégration}

Les tests fonctionnels couvrent les scénarios visibles par l'utilisateur. Les tests d'intégration vérifient le fonctionnement conjoint des composants : formulaire Angular, service Angular, API, service backend, repository et base SQL Server. Cette distinction permet de détecter aussi bien les erreurs d'interface que les erreurs de contrat ou de persistance.

\subsection{Optimisation des performances}

Les optimisations portent principalement sur le chargement Angular, la limitation des données retournées, la pagination ou limitation des listes et la réduction des calculs inutiles lors du rendu documentaire. Les vues de tables utilisent par exemple une limite de lignes afin d'éviter de charger des volumes trop importants.

\subsection{Préparation de la livraison}

La préparation de la livraison inclut la documentation, les captures d'écran, les diagrammes, la vérification des prérequis techniques et la préparation d'un scénario de démonstration. L'objectif est de présenter une application cohérente et compréhensible par un jury.

\section{Analyse et conception}

\subsection{Diagramme de cas d'utilisation}

\diagramplaceholder{sprint5_cas_utilisation}{Diagramme de cas d'utilisation du sprint 5}{Cas d'utilisation d'intégration, qualité et livraison}

Le diagramme de cas d'utilisation du sprint 5 présente les activités transversales : valider un parcours complet, exécuter les tests, vérifier la sécurité, préparer la documentation et réaliser la démonstration. Les acteurs sont l'utilisateur final, l'administrateur, le super administrateur, le développeur et l'évaluateur PFE.

Ce diagramme montre que la livraison ne concerne pas uniquement le développeur. Elle implique aussi les utilisateurs qui valident les parcours, l'administrateur qui contrôle les données métier et l'évaluateur qui doit comprendre l'architecture et reproduire les scénarios.

\subsection{Diagramme de classes}

\diagramplaceholder{sprint5_classes}{Diagramme de classes du sprint 5}{Classes et services transversaux de la phase d'intégration}

Le diagramme de classes du sprint 5 met en évidence les composants transversaux : \texttt{ApiService}, \texttt{AuthInterceptor}, \texttt{ErrorInterceptor}, \texttt{EditorStateService}, \texttt{TemplateService}, \texttt{FamilyService}, \texttt{TableViewService} et \texttt{DocumentRenderService}. Ces services constituent la colonne vertébrale de l'application Angular.

Côté backend, \texttt{EditorController}, \texttt{AuthController}, \texttt{EditorService} et \texttt{EditorRepository} assurent l'orchestration des fonctionnalités. La phase d'intégration vérifie que ces composants collaborent correctement sur les parcours complets.

\subsection{Diagrammes de séquence}

\diagramplaceholder{sprint5_sequences}{Diagramme de séquence du sprint 5}{Séquence complète de démonstration finale}

Le diagramme de séquence décrit un parcours de démonstration complet. L'utilisateur se connecte, le frontend récupère l'état applicatif, l'administrateur vérifie les modèles et chartes, l'utilisateur sélectionne une famille, choisit un bénéficiaire, génère une prévisualisation, puis consulte le résultat. Chaque étape traverse plusieurs couches et constitue donc un test d'intégration réel.

\subsubsection*{Description des interactions système}

Le sprint 5 vérifie la continuité entre les modules développés. Un problème dans l'authentification empêcherait le chargement de l'état. Une erreur dans les familles empêcherait la sélection du bénéficiaire. Une incohérence dans les modèles provoquerait un rendu incomplet. La phase d'intégration consiste donc à tester les dépendances fonctionnelles dans leur ordre réel d'utilisation.

\subsubsection*{Analyse technique}

L'analyse technique porte sur la robustesse des appels HTTP, la gestion des erreurs, la sécurité des endpoints et la performance des opérations. Les interceptors Angular jouent un rôle central : l'interceptor d'authentification ajoute le token, l'interceptor d'erreur peut intercepter les réponses non autorisées et déclencher une déconnexion ou un message utilisateur.

\subsubsection*{Architecture des composants concernés}

L'architecture finale repose sur une communication REST entre Angular et ASP.NET Core. Angular fournit les pages et services de présentation, tandis que l'API centralise les règles métier et l'accès aux données. SQL Server conserve les utilisateurs, les configurations, les modèles et les données métier. Cette architecture facilite la maintenance, car chaque couche dispose de responsabilités identifiables.

\section{Réalisation}

\subsection{Validation des parcours utilisateurs}

Les parcours utilisateurs validés couvrent les trois rôles. Le super administrateur se connecte, accède à l'espace global, configure les familles et les vues de tables. L'administrateur gère les modèles, modifie les contenus et associe les chartes graphiques. L'utilisateur final consulte les familles disponibles, sélectionne un modèle, choisit un bénéficiaire et prévisualise un document.

\screenshotplaceholder{sprint5_demo.png}{Parcours de démonstration finale}{Démonstration d'un parcours complet de génération documentaire}

La validation a également porté sur les redirections automatiques. Un utilisateur déjà connecté ne doit pas revenir sur la page de connexion. Un utilisateur non autorisé ne doit pas accéder à un espace qui ne correspond pas à son rôle.

\subsection{Gestion des erreurs}

La gestion des erreurs est indispensable pour éviter les blocages silencieux. Côté Angular, les interceptors permettent de centraliser une partie du traitement. Côté backend, les contrôleurs retournent des réponses structurées lorsque des paramètres obligatoires sont absents ou lorsqu'une ressource n'existe pas.

\begin{lstlisting}[caption={Principe d'un interceptor d'erreur Angular}]
export const errorInterceptor: HttpInterceptorFn = (req, next) => {
  const authService = inject(AuthService);
  return next(req).pipe(
    catchError((error) => {
      if (error.status === 401) {
        authService.logout();
      }
      return throwError(() => error);
    })
  );
};
\end{lstlisting}

Ce type de mécanisme améliore l'expérience utilisateur et évite qu'une session expirée produise des erreurs répétées dans les pages métier.

\subsection{Optimisation et renforcement de la sécurité}

Le renforcement de la sécurité s'est appuyé sur plusieurs vérifications. Les endpoints critiques doivent être protégés par JWT. Les données envoyées par le client doivent être validées. Les requêtes SQL dynamiques doivent limiter les champs manipulés et utiliser des paramètres. Les mots de passe ne doivent jamais être exposés dans les réponses HTTP.

Les optimisations portent également sur le chargement côté frontend. Les routes utilisent le chargement paresseux des composants, ce qui évite de charger immédiatement toutes les pages. Les services d'état évitent des appels API répétés en conservant les ressources déjà chargées.

\subsection{Documentation technique et utilisateur}

La documentation technique présente l'architecture générale, les endpoints API, les entités principales, les choix technologiques et les diagrammes UML. La documentation utilisateur décrit les principales actions : connexion, navigation selon le rôle, création d'une famille, création d'un modèle, configuration d'une charte graphique, configuration d'une vue de table et prévisualisation d'un document.

Les captures d'écran doivent être insérées dans le dossier \texttt{images/}. Les diagrammes Draw.io doivent être conservés dans le dossier \texttt{diagrammes/}, puis exportés en PNG pour être inclus dans le rapport final.

\subsection{Préparation de la démonstration PFE}

La démonstration PFE doit suivre un scénario clair et progressif. Elle commence par la connexion d'un super administrateur, se poursuit par la présentation des modules d'administration, puis montre l'édition d'un modèle et se termine par la prévisualisation d'un document généré à partir d'un bénéficiaire.

Le scénario recommandé est le suivant :

\begin{enumerate}
    \item connexion avec un compte de démonstration ;
    \item présentation rapide du tableau de bord ou de l'espace correspondant au rôle ;
    \item création ou consultation d'une famille de documents ;
    \item consultation d'un modèle associé à cette famille ;
    \item modification mineure du contenu dans l'éditeur Tiptap ;
    \item association ou visualisation d'une charte graphique ;
    \item sélection d'un bénéficiaire ;
    \item génération de la prévisualisation ;
    \item explication de l'architecture backend/frontend et des contrôles de sécurité.
\end{enumerate}

\screenshotplaceholder{sprint5_tests.png}{Résultats de tests et validation technique}{Résultat des tests, build ou validation finale}

\section{Tests et validation}

\begin{longtable}{p{0.25\textwidth} p{0.43\textwidth} p{0.22\textwidth}}
\caption{Scénarios de test du sprint 5}\\
\toprule
\textbf{Scénario} & \textbf{Description} & \textbf{Résultat attendu}\\
\midrule
\endfirsthead
\toprule
\textbf{Scénario} & \textbf{Description} & \textbf{Résultat attendu}\\
\midrule
\endhead
Parcours super administrateur & Connexion, accès espace global, gestion famille et vue table. & Parcours terminé sans erreur bloquante.\\
\midrule
Parcours administrateur & Connexion, modification modèle, association charte. & Modèle sauvegardé et rendu mis à jour.\\
\midrule
Parcours utilisateur final & Connexion, choix famille, choix bénéficiaire, prévisualisation. & Document prévisualisé correctement.\\
\midrule
Exploitation module métier & Accès à l'espace données, sélection d'un module et bascule entre les tables. & Affichage cohérent des tables regroupées.\\
\midrule
Expiration session & Utilisation d'un token expiré ou supprimé. & Redirection vers la connexion.\\
\midrule
CRUD complet & Création, modification et suppression de ressources métier. & Données persistées et état frontend cohérent.\\
\midrule
Build frontend & Compilation Angular en mode production. & Build terminé sans erreur critique.\\
\midrule
Build backend & Compilation du projet ASP.NET Core. & Build terminé sans erreur critique.\\
\bottomrule
\end{longtable}

Les tests de sécurité ont vérifié le refus des appels anonymes sur les endpoints protégés, l'absence du mot de passe haché dans les réponses et la redirection des utilisateurs vers leur espace autorisé. Les tests d'intégration ont validé les flux entre Angular, l'API et SQL Server. Les tests fonctionnels ont confirmé que les principaux scénarios de démonstration étaient réalisables.

\section{Conclusion}

Le sprint 5 a permis de consolider l'ensemble du projet et de préparer sa livraison. Les modules développés dans les sprints précédents ont été intégrés dans des parcours cohérents. Les tests ont permis d'identifier les points critiques, de vérifier les contrats API et de valider les scénarios nécessaires à la démonstration.

À l'issue de ce sprint, la plateforme dispose d'une base fonctionnelle complète : authentification, rôles, édition documentaire, administration métier, configuration des vues de tables, gestion des bénéficiaires et prévisualisation dynamique. Le projet est ainsi prêt à être présenté comme une solution professionnelle de gestion documentaire dynamique.


%=================================================
% CONCLUSION
%=================================================

\chapter*{Conclusion générale et perspectives}
\addcontentsline{toc}{chapter}{Conclusion générale et perspectives}

\section*{Perspectives}
\addcontentsline{toc}{section}{Perspectives}

Pour la suite du projet, plusieurs perspectives d'évolution peuvent être envisagées afin d'enrichir SIRH-Doc et de l'adapter à des usages administratifs plus larges. Une première évolution concernerait l'intégration d'un circuit de validation documentaire. Dans ce scénario, un document pourrait être préparé par un utilisateur, vérifié par un responsable, puis validé avant son archivage définitif. Cette fonctionnalité renforcerait la fiabilité des pièces produites, surtout pour les documents sensibles comme les contrats, les attestations ou les décisions administratives.

Une autre perspective importante serait l'ajout de la signature électronique. Elle permettrait de donner une valeur plus forte aux documents générés et de réduire davantage les échanges papier. Cette évolution pourrait être complétée par l'envoi automatique des documents par courriel, notamment pour les vacataires, les enseignants ou les agents qui doivent recevoir rapidement une copie officielle.


\section*{Conclusion générale}
\addcontentsline{toc}{section}{Conclusion générale}

La réalisation de SIRH-Doc a constitué une expérience complète d'ingénierie logicielle. Le projet ne s'est pas limité à la création d'une interface de génération de documents ; il a demandé de comprendre un contexte métier réel, de tenir compte des contraintes d'un établissement universitaire et de concevoir une solution capable de rester adaptable d'une organisation à l'autre.

Tout au long du projet, plusieurs connaissances acquises durant la formation ont été mobilisées : la modélisation UML, la conception orientée objet, l'architecture logicielle, le développement web, la gestion des bases de données et l'organisation du travail par sprints. Ces éléments ont été appliqués dans un cadre concret, avec des choix techniques qui devaient répondre à des besoins précis : isoler les données, gérer les rôles, prendre en compte l'année universitaire, configurer les modules et assurer la traçabilité des documents produits. Ce travail nous a aussi permis de mieux comprendre la manière dont un projet est conduit dans un environnement professionnel, avec ses contraintes, ses échanges et ses ajustements progressifs.

Le projet a également permis d'approfondir de nouvelles compétences. La mise en place d'un moteur basé sur les métadonnées, la génération documentaire à partir de templates, la gestion des variables tableau et l'archivage des documents ont posé des problèmes qui dépassent les exercices classiques de développement. L'utilisation de Tiptap nous a notamment fait découvrir plus concrètement le monde de l'open source et la manière d'adapter une bibliothèque existante à un besoin métier précis.

Malgré les contraintes de temps liées au cadre d'un projet de fin d'études, le travail réalisé a permis d'aboutir à une plateforme cohérente et fonctionnelle. Les quatre sprints ont progressivement construit le socle de sécurité, le moteur de données, le studio documentaire puis le module de production et d'archivage. Cette progression a donné au projet une structure lisible et a facilité la validation des fonctionnalités au fur et à mesure.

En définitive, SIRH-Doc représente une base solide pour accompagner la dématérialisation des documents RH dans un établissement éducatif. Le système reste perfectible, comme tout produit appelé à évoluer, mais ses fondations techniques et fonctionnelles permettent déjà d'envisager des extensions sans remettre en cause l'ensemble de l'architecture. Ce projet a ainsi été l'occasion de relier les acquis académiques à une problématique réelle, tout en développant une solution utile, structurée et ouverte à de futures améliorations.


%=================================================
% BIBLIOGRAPHIE ET WEBOGRAPHIE
%=================================================

\begin{thebibliography}{9}
\addcontentsline{toc}{chapter}{Bibliographie}
\bibitem{pressman}
R. S. Pressman et B. R. Maxim, \textit{Software Engineering: A Practitioner's Approach}, McGraw-Hill, 2020.
\bibitem{uml}
Object Management Group, \textit{Unified Modeling Language Specification}, OMG.
\bibitem{scrum}
K. Schwaber et J. Sutherland, \textit{The Scrum Guide}, 2020.
\end{thebibliography}

\chapter*{Webographie}
\addcontentsline{toc}{chapter}{Webographie}
\begin{itemize}
    \item Documentation Microsoft ASP.NET Core : \url{https://learn.microsoft.com/aspnet/core}
    \item Documentation Angular : \url{https://angular.dev}
    \item Documentation SQL Server : \url{https://learn.microsoft.com/sql/sql-server}
    \item Documentation Tiptap : \url{https://tiptap.dev}
    \item Documentation Draw.io : \url{https://www.drawio.com}
\end{itemize}

%=================================================
% ANNEXES
%=================================================

\chapter*{Annexes}
\addcontentsline{toc}{chapter}{Annexes}

\section*{Annexe A : Organisation des diagrammes UML}

Les diagrammes UML du projet sont regroupés dans le dossier \texttt{diagrammes/}. Chaque fichier de sprint contient plusieurs pages Draw.io afin de garder dans un même livrable les vues de conception d'une même itération.

\begin{longtable}{p{4cm} p{9cm}}
\toprule
\textbf{Fichier} & \textbf{Contenu} \\
\midrule
\texttt{classe\_global.drawio} & Diagramme de classes global et diagramme de cas d'utilisation global. \\
\texttt{sprint1.drawio} & Classes, cas d'utilisation, séquence d'authentification JWT et séquence de résolution tenant. \\
\texttt{sprint2.drawio} & Classes, cas d'utilisation, séquence de lecture dynamique et séquence d'importation intelligente. \\
\texttt{sprint3.drawio} & Classes, cas d'utilisation, séquence de sauvegarde de template et séquence de prévisualisation avec fusion. \\
\texttt{sprint4.drawio} & Classes, cas d'utilisation, séquence de génération en masse et séquence d'audit soft-delete. \\
\bottomrule
\end{longtable}

\section*{Annexe B : Captures d'écran à intégrer}

Les captures suivantes peuvent être ajoutées dans le dossier \texttt{images/} au moment de la préparation de la version finale imprimable :

\begin{itemize}
    \item interface de connexion et redirection selon le rôle ;
    \item écran de sélection de l'année universitaire ;
    \item tableau de bord administrateur ;
    \item configuration d'un module métier et d'une TableView ;
    \item configuration d'une famille documentaire ;
    \item studio d'édition d'un template ;
    \item prévisualisation d'un document fusionné ;
    \item écran de génération documentaire ;
    \item liste des archives avec filtres ;
    \item détail d'une archive et affichage du document généré.
\end{itemize}

\section*{Annexe C : Exemples de composants techniques}

Les éléments suivants constituent des points d'entrée utiles pour relier le rapport au code source :

\begin{itemize}
    \item \texttt{backend/Middlewares/TenantResolutionMiddleware.cs} pour la résolution du tenant ;
    \item \texttt{backend/Common/Tenant/TenantResolver.cs} pour la lecture de l'organisation ;
    \item \texttt{backend/Repositories/EditorRepository.cs} pour l'accès SQL dynamique ;
    \item \texttt{frontend/src/app/features/editor/services/document-render.service.ts} pour le rendu documentaire ;
    \item \texttt{frontend/src/app/features/editor/services/table-view.service.ts} pour la gestion des vues dynamiques ;
    \item \texttt{frontend/src/app/features/editor/pages/document-archive/} pour la consultation des archives.
\end{itemize}


\end{document}
